% concatenated from InvariantTightContactStructures.tex
\documentclass[a4paper, 11pt]{amsart}
\usepackage[T1]{fontenc}
\usepackage[utf8]{inputenc}
\usepackage[english]{babel}

\usepackage{microtype}

\usepackage{mathtools, amssymb}

\usepackage{tikz-cd}
\usetikzlibrary{arrows.meta, bending}
\usetikzlibrary{decorations.markings}

\usepackage{enumitem}
\setlist[enumerate, 1]{label={\arabic*)}}

\usepackage[autostyle]{csquotes}
\usepackage[backend=biber, style=alphabetic, hyperref, giveninits=true]{biblatex}
\addbibresource{InvariantTightContactStructures.bib}

\usepackage{refcount}
\usepackage{graphicx}

\usepackage{hyperref}
\hypersetup{hidelinks}
\usepackage{aliascnt}
\usepackage{cleveref}






\theoremstyle{plain}
\newtheorem{prop}{Proposition}[section]
\Crefname{prop}{Proposition}{Propositions}
\newaliascnt{cor}{prop}
\newtheorem{cor}[cor]{Corollary}
\aliascntresetthe{cor}
\newaliascnt{lemma}{prop}
\newtheorem{lemma}[lemma]{Lemma}
\aliascntresetthe{lemma}
\newtheorem{thm}{Theorem}
\Crefname{thm}{Theorem}{Theorems}
\newtheorem{thmR}{Theorem}

\theoremstyle{definition}
\newaliascnt{defin}{prop}
\newtheorem{defin}[defin]{Definition}
\aliascntresetthe{defin}
\newtheorem*{definsn}{Definition}
\newaliascnt{rmk}{prop}
\newtheorem{rmk}[rmk]{Remark}
\aliascntresetthe{rmk}

\newcommand{\R}{\mathbb{R}}
\newcommand{\sd}{\text{std}}
\newcommand{\XX}{\mathfrak{X}}
\newcommand{\Lie}{\mathcal{L}}
\newcommand{\FF}{\mathcal{F}}
\newcommand{\NN}{\mathcal{N}}
\newcommand{\inv}{\text{inv}}
\newcommand{\zero}{\mathbf{0}}
\newcommand{\id}{\text{id}}
\newcommand{\Z}{\mathbb{Z}}
\newcommand{\emb}{\text{emb}}
\newcommand{\bimath}{\bar{\imath}}
\renewcommand{\epsilon}{\varepsilon}
\renewcommand{\theta}{\vartheta}
\DeclareMathOperator{\Emb}{Emb}
\DeclareMathOperator{\Loop}{Loop}
\DeclareMathOperator{\tb}{tb}
\DeclareMathOperator{\Diff}{Diff}
\DeclareMathOperator{\dist}{dist}
\DeclareMathOperator{\Tight}{Tight}
\DeclareMathOperator{\Imm}{Im}
\DeclareMathOperator{\Ker}{Ker}
\DeclareMathOperator{\diver}{div}
\DeclareMathOperator{\Tr}{Tr}
\DeclarePairedDelimiter{\ang}{\langle}{\rangle}
\DeclarePairedDelimiter{\abs}{\lvert}{\rvert}
\DeclarePairedDelimiter{\Abs}{\lVert}{\rVert}


\begin{document}
\title[Invariant tight contact structures]{Classification of invariant tight contact structures on the 3-space, -ball and -sphere}
\author{Mirko Torresani}
\address{DIMAI, viale Giovan Battista Morgagni 67/a, 50134, Firenze, Italy}
\subjclass[2020]{Primary 57K33, 57M60; Secondary 57K10}
\begin{abstract}
	We prove some classification results for tight contact structure in the 3-space, -ball and -sphere that are invariant with respect to some arbitrary involution, that is conjugated to the standard rotation around the x-axis. Unlike the classical scenario, a new integral torsion appears, dictating a splitting between equivalence classes. These tools could be useful fur future classification results regarding strongly invertible Legendrian knots.
\end{abstract}
\maketitle
\tableofcontents
\section{Introduction}\label{sec:introdution}
In recent year the study of strongly invertible Legendrian knots has received considerable attention. The aim of this paper is to build some basic classification results for tight contact structures in $\R^3$, $S^3$ and $B^3_{\sd}$ that are invariant with respect to some $\Z/2\Z$-action on such spaces. 


Recall that a contact structure $\xi$ on an oriented three manifold $M$ is a plane distribution such that, if locally we fix a form $\alpha$ whose Kernel is $\xi$, then $\alpha \wedge d\alpha$ is a non-vanishing 3-form. The structure $\xi$ is \emph{positively cooriented} if it coincides with $\Ker(\alpha)$ for some global 1-form $\alpha$, such that $\alpha \wedge d\alpha$ is positive with respect of the orientation on $M$. From now we will always suppose such condition to hold.

A contact structure $\xi$ is \emph{overtwisted} if there exists an embedded disk $D$ with $T_pD$ intersecting $\xi_p$ only at $T_p\partial D$ for every $p \in \partial D$. If $\xi$ is not overtwisted, then it is \emph{tight}.

Furthermore, an oriented surface $S$, in a neighborhood $S \times \R$ the contact form can be written as $\beta_z + w_z\,dz$, with $w_z \colon S \to \R$ a smooth function and $\beta_z$ identifiable a 1-form $S$. The contact condition is given by
\begin{equation}\label{eq:contcond}
	w_z \,d\beta_z +\beta_z \wedge(dw_z - \dot{\beta}_z) > 0\,.
\end{equation}
Moreover, if we equip $S$ of an area form $\Omega$, then we define the \emph{characteristic foliation} $S_\xi$ to be the unique vector field $X$ such that $i_X\Omega = \beta_0$, where $i_X$ is the contraction given by substituting $X$ in the \emph{first entry}. The singularity of $X$ are the points where the planes $\xi_p$ and $T_pS$ coincide. we give signs to such singularities: \emph{positive} if such planes have same orientation, \emph{negative} otherwise.

Finally, given a global contact form $\alpha$, the \emph{Reeb vector field} is the unique vector field $R_\alpha \in \mathfrak{X}(\R^3)$ such that 
\[
\begin{split}
	d\alpha(R_\alpha,\cdot) &\equiv 0 \\
	\alpha(R_\alpha) &\equiv 1.
\end{split}
\]

Let's now introduce the involution we will be working with. A priori we want to fix a generic involution $u$, acting on $\R^3$, with the following properties:
\begin{enumerate}
	\item it extends to an infinity-preserving involution of $S^3$;
	\item the fixed point set is non-empty, at every fixed point the differential, viewed as an endomorphism of the tangent space, is diagonalizable as $[1] \oplus [-1] \oplus [-1]$.
\end{enumerate}
Thanks to Waldhausen's results, such $u$ is conjugated to a standard rotation. However, it will be more useful to abstract the reasoning up to the last section.

We say that a contact structure $\xi$ is \emph{invariant} if $u^*\xi = \xi$ and $u^*\alpha = -\alpha$. The aim is to classify such objects up to \emph{equivariant isotopy}, i.e.\ up to isotopies $\psi_t$ of the space such that $\psi_t \circ u$ coincides with $u \circ \psi_t$. The three main theorems are then the following. To future reference our isotopies will always start at the identity.

\begin{thm}\label{thmmainthmRthree}
	The set of tight contact structures on $\R^3$ up to equivariant isotopy is trivial.
\end{thm}
\begin{thm}\label{thmmainthmDthreecomplementary}
	If such action $u$ preserves the complementary of the standard $3$-ball, then the set of invariant tight contact structures on $(B^3_{\sd})^c$, inducing a fixed characteristic foliation on $S^3 = \partial B^3_{\sd}$, up to equivariant isotopy is trivial.
\end{thm}
\begin{thm}\label{thmmainthmDthree}
	If such action $u$ preserves the standard $3$-ball, then the set of invariant tight contact structures on $B^3_{\sd}$, inducing a fixed characteristic foliation on $S^3 = \partial B^3_{\sd}$, up to equivariant isotopy is a $\Z$-torsor.
\end{thm}
\begin{thm}\label{thmmainthmSthree}
	The set of invariant tight contact structures on $S^3$, up to equivariant isotopy, is a $\Z$-torsor.
\end{thm}

Before passing to the next section, let's briefly define the surfaces, and vector fields, we will work with. All the surfaces will be oriented and connected.
\begin{definsn}
	A surface $S \subseteq \R^3$ preserved by $u$ is said to be $u$-oriented if it is equipped with an area form $\Omega \in \Lambda^2S$ such that $u^*\Omega = \Omega$.
\end{definsn}
\begin{definsn}
	Given an $u$-oriented surface $S$, an anti-invariant singular foliation on $S$ is an element of
	\[
		\{X \in \XX(S) \mid u^*X = -X\}/\sim\,,
	\]
	where
	\begin{enumerate}
		\item $u^*X$ at a point $p$ is defined as the pull-back, via the differential, of $X(u(p))$;
		\item $X$ is equivalent to $X'$ if there exists a smooth function $f \colon Z \to \R^+$ such that $f \circ u = f$ and $X' = f\,X$.
	\end{enumerate}
\end{definsn}
The equivalence class of the characteristic foliation on a $u$-oriented surface induced by an invariant contact structure is an example of such foliation.

For now on the term \emph{anti-invariant} will be omitted, and we will refer to the previous vector fields simply as \emph{singular foliation}.


The paper will be structured as follows. In \Cref{sec:EThm} we state and prove equivariant versions of classical theorems in contact topology; \Cref{sec:EConvex} is dedicated to expand on this theorem in order to set foundation for a basic equivariant convex theory. Although the setting is new, most of the work consists in checking that the $u$-action respects the quantity at play. The last preparations are done in \Cref{sec:tomography}, in which a classification theorem for invariant tight contact structure in $S^2 \times [-1,1]$ is proven. Here is where the equivariant scenario diverges from the classical one: a new $\Z$-torsor appears, dictating a splitting between equivalence classes. Finally \Cref{sec:proofmainthmopen,sec:proofmainthmclose} are devoted to proving the main results.
\section{Equivariant Giroux and Darboux' Theorems}\label{sec:EThm}
The aim of this section is to state and prove equivariant versions of classical results in contact topology: Gray, Giroux and Darboux' theorems. Such restatements are given by \Cref{EGray,EGiroux,EDarbouxustd} respectively.

Let's introduce the following notation: given a surface $S$, we will extend the action on $S \times \R$ as $u(x,y) = (u(x), y)$.
\begin{prop}\label{tubsurf}
	Each $u$-oriented surface admits a tubular neighborhood diffeomorphic $S \times \R$, which is $u$-equivariant, and such that the action is given by the convention above.
\end{prop}
\begin{proof}
	Since $u$ is an involution, it's action can be viewed as a $\Z/2\Z$, which is a compact group. Hence we can fix an invariant Riemannian metric $g$ on $S$. By Bredon \cite[Thm.\ 2.2]{bredonIntroductonCompactTransformation1972} there exists small enough invariant neighborhood $U$ of $S$ inside the normal bundle $N(S) = TS^\bot \subseteq T\R^3$ such that the exponential map $\exp \colon U \to M$ is an embedding. Let's observe that the exponential map is equivariant, with the action on the normal bundle given by the differential.
	
	Then the image $\exp(U)$ is the tubular neighborhood we searched for. Moreover, since $N(S)$ is an oriented 1-bundle, it is isomorphic to $S \times \R$ as vector bundle. By taking at each fiber the positively oriented vector of $g$-norm 1, one can impose a canonical isomorphism.
	
	Finally, since $u$ is a $g$-isometry, it must send $(x,y)$ into $(u(x),\pm y)$. Since $u$ preserves the orientation both of $S$ and $\R^3$ such sign must be negative.
\end{proof}

\begin{prop}\label{EGray}
	Let $\xi_t$, $t \in [0,1]$, be a smooth family of invariant contact structures on $S^3$. Then there exist an equivariant isotopy $\psi_t$ such that $T\psi_t(\xi_0) = \xi_t$. 
\end{prop}
\begin{proof}
	The equation we have to solve is
	\[
		\psi_t^*\alpha_t = \lambda_t\alpha_0
	\]
	for some family of functions $\lambda_t \colon S^3 \to \R$. By taking derivatives in the previous equation we get
	\[
		\psi^*_t (\dot{\alpha}_t + \Lie_{X_t}\alpha_t) = \dot{\lambda}_t\alpha_0 = \frac{\dot{\lambda}_t}{\lambda_t}\psi_t^*\alpha_t\,,
	\]
	where the dot denotes the derivative with respect to $t$ \cite[Lemma 2.2.1]{geigesIntroductionContactTopology2008}.
	
	Let's set $\mu_t \coloneqq \tfrac{d}{dt}(\log \lambda_t )\circ \psi_t^{-1}$, and let's use the Cartan's formula\footnote{which holds also for time-dependent vector field, see \cite[App.\ B]{geigesIntroductionContactTopology2008}.} to get
	\[
		\psi^*_t (\dot{\alpha}_t + d(\alpha_t(X_t)) +\iota_{X_t}\alpha_t) = \dot{\lambda}_t\alpha_0 = \psi_t^*(\mu_t\alpha_t)\,.
	\]
	
	By making the Ansatz $X_t \in \xi_t$, the previous equation implies
	\begin{equation}\label{eq:Xt}
	\begin{cases}
		\dot{\alpha}_t + i_{X_t}d\alpha_t = \mu_t\alpha_t \\
		\dot{\alpha}_t(R_t) = \mu_t
	\end{cases}\,,
	\end{equation}
	where $R_t$ is the Reeb's vector field of $\alpha_t$.
	
	Finally, having defined $\mu_t$, we can choose a unique element $X_t \in \xi_t$ satisfying Equations \eqref{eq:Xt}: restricting the 1-forms to $\xi_t$ allows us to pinpoint a precise $X_t$, thanks to the non-degeneracy of $d\alpha_t$, and with such choice Equations \eqref{eq:Xt} holds, since $R_t \in \Ker(\mu_t\alpha_t - \dot{\alpha}_t)$.
	
	Moreover
	\[
		\mu_t \circ u = u^*\mu_t = u^*[\dot{\alpha}_t(R_t)] = - \dot{\alpha}_t(- R_t) = \mu_t\,.
	\] 
	
	Hence $u^*X_t$ satisfies the following equality:
	\[
		u^*\dot{\alpha}_t + i_{u^*X_t}du^*\alpha_t = \mu_t u^*\alpha_t\,,
	\]
	re-writable as
	\[
		-\dot{\alpha}_t - i_{u^*X_t} d\alpha_t = -\mu_t\alpha_t\,.
	\]
	Thus $u^*X_t$ satisfies the same equation as $X_t$, and hence they must coincide. By integrating $X_t$ we obtain an equivariant isotopy $\psi_t$ with the desired properties.
\end{proof}
The previous discussion also holds for contact structures in $\R^3$. However, the last step is no longer true, and we can integrate $X_t$ on $[0,1]$ only over a compact set, or on the whole $\R^3$ if $X_t$ vanishes away of such set.
\begin{prop}\label{EGiroux}
	Let $S_0$ and $S_1$ two  $u$-oriented compact surfaces in $\R^3$, and let $\phi \colon S_0 \to S_1$ an equivariant diffeomorphism, such that $\phi(S_{0,\xi_0}) = S_{1,\xi_1}$ singular foliations, for two invariant contact structures. Then there exist appropriate invariant neighborhoods and an equivariant contactomorphism $\psi \colon \NN(S_0) \to \NN(S_1)$ such that $\psi|_{S_0}$ is isotopic to $\psi$ via an isotopy preserving the characteristic foliation.
\end{prop}
\begin{proof}
	On invariant tubular neighborhoods of $S_0$ and $S_1$, diffeomorphic to $S_i \times \R$, the contact forms can be written as $\beta_z^i + w_z^idz$, with $\beta_z^i$ and $u_z^i$ that are anti-invariant.
	
	By hypothesis there exists a function  $f \colon S \to \R$ such that $\beta_0^0$ coincides with $f \phi^*\beta_0^1$, with $f$ being constantly positive or negative. The sign depends on the fact that $\phi$ does or does not respects the orientations of $S_0$ and $S_1$. Moreover, $f$ is invariant with respect to the the $u$-action. 
	
	Secondly, let's consider the natural extension of $\phi$ from $S_0 \times \R$ to $S_1 \times \R$, which we identify with invariant neighborhoods, such that the contact form $f\,\phi^*\alpha_1$ is positively cooriented in $S_0 \times \R$.
	
	Let's define the family of forms $\alpha_t$ on $\NN(S_0) \cong  S_0 \times \R$ as
	\[
		\alpha_t \coloneqq (1-t)\alpha_0 + tf \phi^*\alpha_1\,.
	\]
	Equation \eqref{eq:contcond} is simultaneously linear in $\dot{\beta}_z$ and $w_z$, and $\alpha_t$ induce the same form $\beta_0$ on $S \times \{0\}$. Hence the contact condition is satisfied on $S \times \{0\}$, and it will be satisfied in a small neighborhood $\abs{z} < 1$. Moreover $\alpha_t$ is always anti-invariant.
	
	By \Cref{EGray} we can construct an equivariant contact isotopy in a small neighborhood $\NN(S_0)$ of $S_0$, keeping $S_0$ fixed. Let's call $X_t$ the vector field $\tfrac{d}{dt}\psi_t$.
	
	On the space $TS_0$ we have that $\dot{\alpha}_t \equiv 0$, since $\alpha_0$ and $f\phi^*\alpha_1$ induce the same form $\beta_0$. Hence, for every $v$ tangent to $(S_0)_{\xi_0}$ we have that
	\[
		\Lie_{X_t}\alpha_t(v) \overset{(\triangle)}{=} i_{X_t}d\alpha_t(v) + di_{X_t}\alpha_t(v) \overset{X_t \in \xi_t}{=} i_{X_t}d\alpha_t(v) \overset{(\ast)}{=} -\dot{\alpha}_t(v) = 0\,,
	\]
	where equality $(\triangle)$ is justified by Cartan's formula, and equality $(\ast)$ by Equations \eqref{eq:Xt}, present in the proof of \Cref{EGray}. As a consequence, the flow of $X_t$ must keep $\beta_0$ fixed, and hence it preserves the characteristic foliation.
	
	Finally, by setting $\psi \coloneqq \phi \circ \psi_1$, the distribution $T\psi(\xi_0)$ coincides with $\xi_1$, and $\psi$ is the required diffeomorphism.
\end{proof}
\begin{prop}\label{Girouxbound}
	With the assumption of the preceding theorem we can find a stronger $\psi$ such that $\psi|_{S_0} = \phi$.
\end{prop}
\begin{proof}
	We want to reuse the proof of \Cref{EGray}, without forcing $X_t$ to belong to $\xi_t$.
	
	As the previous theorem, we know that $\alpha_0$ and $f\phi^*\alpha_1$ induce the same form $\beta_0$ on $S_0 = S_0 \times \{0\}$ for some function $f$; in particular $\dot{\alpha}_t|_{TS_0} \equiv 0$. By no longer requiring that $X_t$ vanishes on $\alpha_t$, the desired condition becomes 
	\[
		\dot{\alpha}_t + d(\alpha_t(X_t)) + i_{X_t}d\alpha_t = \mu_t\alpha_t\,,
	\]
	where
	\[
	\mu_t = \frac{d}{dt}\,\log \lambda_t \circ \psi_t^{-1}\,.
	\]
	
	Let's write $X_t = H_tR_t + Y_t$, with $Y_t \in \Ker \alpha_t$. Then the previous condition becomes 
	\begin{equation}\label{eq:condX}
		\dot{\alpha}_t + dH_t + i_{Y_t} d\alpha_t = \mu_t\alpha_t\,.
	\end{equation}
	
	The Equation \eqref{eq:condX} completely determines $\mu_t$ via $H_t$ once we evaluate it on the Reeb vector field. Once this is done, $Y_t$ is determined by the non-degeneracy $d\alpha_t|_{\xi_t}$.
	
	The function $H_t$ is determined via the following system, whose equations are compatible thanks to the condition $\dot{\alpha}_t|_{TS_0} \equiv 0$:
	\begin{equation}\label{eq:system}
		\begin{cases}
			\dot{\alpha}_t + dH_t = 0 & \text{on $S_0$}\\
			H_t \equiv 0 & \text{on $S_0$}
		\end{cases}\,.
	\end{equation}
	
	At this point the vector field $X_t$ identifies a flow $\psi_t$ fixing $S_0$ pointwise.
	
	Moreover, let's show that $u^*X_t = X_t$ holds. As a start, from Equations \eqref{eq:system} we get 
	\[
	u^*dH_t = -u^*\dot{\alpha}_t =  \dot{\alpha}_t =- dH_t\,.
	\]
	As a consequence, Equation \eqref{eq:condX} implies that $u^*Y_t = Y_t$. Hence $u^*$ preserves also $X_t = H_tR_t + Y_t$.
\end{proof}
\begin{prop}\label{EDarbouxustd}
	Suppose that $u$ is given by $(x,y,z) \mapsto (x,-y-z)$. Given an invariant contact structure in $\R^3$ (or $S^3$), we can find a global equivariant isotopy such that $\xi$ coincides with $\Ker(dz - y\,dx)$ in a small neighborhood of the origin.
\end{prop}
\begin{proof}
	We will build a deformation of the contact structure, and realize it via \Cref{EGray}.
	
	Firstly, let's fix $v \in T_\zero\R^3$ such that $\alpha(\zero)$ coincides with $\ang{v,\cdot }$. Since $u^*\alpha$ equals $-\alpha$, then $d_pu(v)$ coincides with $-v$. 
	
	Thus the vector $v$ belongs to the $yz$-plane. Moreover, up an isotopy of the form
	\[
	\begin{bmatrix}
		1 & 0 & 0 \\
		0 & \cos t  & -\sin t\\
		0 & \sin t & \cos t
	\end{bmatrix}
	\]
	we can suppose $v$ to be a positive multiple of the vector $\partial_z$. Consequently, $\xi(\zero)$ is set to the $xy$-plane. Moreover, without loss of generality, we can suppose that we have chosen the contact form such that $v$ is the vector $\partial_z$. Moreover, $d\alpha(\zero)$ is a symplectic 2-form, positively oriented, on $\xi(\zero)$. Hence $d\alpha|_{\xi(\zero)}$ is equal to $\lambda\,dx \wedge dy$ for some positive $\lambda$.
	
	Thus, up to rescale the $x$- and $y$-axis with an isotopy of the form
	\[
	\psi_t \colon
	\begin{bmatrix}
		x \\
		y \\
		z
	\end{bmatrix}
	\mapsto (1-t)
	\begin{bmatrix}
		x\\[0.5ex]
		y\\[0.5ex]
		z
	\end{bmatrix}
	+ t
	\begin{bmatrix}
		\sqrt{\lambda}\,x\  \\[0.5ex]
		\sqrt{\lambda}\,y \\[0.5ex]
		z
	\end{bmatrix} 
	\]
	and substitute $\xi$ with $T\psi_1^{-1}(\xi)$, we can actually suppose that $d\alpha|_{\xi(\zero)}$ coincides with $dx \wedge dy$.
	
	Now we want to adjust the two-from $d\alpha$ on the whole tangent space $T_\zero\R^3$. Firstly, since it is anti-invariant the equality
	\[
	d\alpha(\zero) = dx\wedge dy + \mu\, dx \wedge dz\,.
	\]
	is forced to hold. Let's then replace $\alpha$ with a new contact form $\alpha'$, for the same structure, given by
	\[
	\alpha_1 \coloneqq e^{-\mu x}\alpha\,.
	\]
	Then 
	\[
	d\alpha_1 = -\mu e^{-\mu x}\,dx \wedge \alpha + e^{-\mu x}\,d\alpha\,,
	\]
	and
	\[
	\begin{cases}
		d\alpha_1(\zero) = dx \wedge dy\\
		\alpha_1(\zero) = \partial_z
	\end{cases}\,.
	\]
	
	
	Finally, just set $\alpha_0 \coloneqq dz - y\,dx$ and interpolate it with $\alpha_1$. Since $\alpha_t$, $d\alpha_t$ at the origin are constantly $\alpha_1$ and $d\alpha_1$ respectively, then close to the origin the interpolation is a contact form. Finally just use a bump function to radially interpolate $\alpha_t$ with $\alpha$ away from the origin. Conclude then via \Cref{EGray}.
\end{proof}
\begin{cor}\label{equalneighorigin}
	Suppose that $u$ is given by $(x,y,z) \mapsto (x,-y,-z)$. Up to equivariant isotopy, any two invariant contact structures on $\R^3$ can be made be equal in a small neighborhood of the origin.
\end{cor}

Let's end the section with a last lemma, that allows us to reconstruct contact structures from singular foliations.
\begin{lemma}\label{charinducescontact}
	Let $S$ be an $u$-oriented compact surface embedded in $\R^3$. A vector field $X$ on $S$ represents the characteristic foliation of some invariant contact structure defined near $S$ if and only if it satisfies the condition
	\[
		\diver_\Omega(X)(p)\neq 0 \text{ if } X(p)=0\,.
	\]
\end{lemma}
\begin{proof}
	Let's prove one implication, by fixing an invariant contact structure $\xi$ near $S$ and a vector field $X$ defining the characteristic foliation. Locally the contact form can be written as $\beta_z + w_zdz$, with $\beta \coloneqq \beta_0$. If $p$ is a singularity of $X$ we know that the planes $T_pS$ and $\xi_p$ coincide. Hence
	\[
		\diver_\Omega(X)(p)\,\Omega_p= (d(i_X\Omega))_p = (d\beta)_p = d\alpha|_{T_pS} = d\alpha|_{\xi_p} \neq 0\,.
	\] 
	Hence $\diver_\Omega(X)$ cannot vanish on any singular point of $X$.
	
	On the other hand, let's fix a vector field $X$ on $S$ satisfying the divergence condition, and set $w \colon S \to \R$ such that $d\beta = w\Omega$, where $\beta \coloneqq i_X\Omega$. By construction $u^*w = - w$.
	
	Observe that the vanishing set of $\beta$ and $X$ coincide. And if $p$ such of a point, then
	\[
		0 \neq \diver_\Omega(X)(p)\,\Omega_p = (d(i_X\Omega))_p = (d\beta)_p\,.
	\]
	Hence $w(p)$ cannot vanish on such points.
	
	Secondly, fix an invariant Riemannian metric $g$ on $S$. Since $S$ is oriented, we know that there is a unique $\gamma \in \Lambda^1S$ such that
	\[
	\begin{cases}
		\Abs{\gamma}_g = \Abs{\beta}_g\\
		\ang{\beta,\gamma}_g = 0\\
		(\beta \wedge \gamma)/\Omega \ge 0\\
		(\beta \wedge \gamma)/\Omega > 0 & \text{if $\beta(p) \neq 0$}
	\end{cases}
	\]
	Moreover, we know that 
	\[
	\begin{cases}
		\Abs{- u^*\gamma}_g = \Abs{\gamma}_g = \Abs{\beta}_g\\
		\ang{\beta, - u^*\gamma}_g = \ang{u^*\beta, u^*\gamma}_g = \ang{\beta,\gamma}_g = 0\\
		(\beta \wedge - u^*\gamma)/\Omega = u^*[(\beta \wedge \gamma)/\Omega] \ge 0
	\end{cases}\,,
	\]
	and hence $u^*\gamma$ coincides with $-\gamma$.
	
	Now set
	\[
		\beta_z \coloneqq \beta + z(dw - \gamma)
	\]
	and
	\[
		\alpha' \coloneqq \beta_z + w\,dz\,.
	\]
	
	Then $u^*\alpha' = - \alpha'$, and
	\[
		w\,d\beta_0 + \beta_0 \wedge (dw - \dot{\beta}_z|_{z = 0}) = w^2\Omega + \beta \wedge \gamma > 0,.
	\]
	Thus, $\alpha'$ defines a invariant contact structure near $S$.
\end{proof}
\section{Equivariant Convex Theory}\label{sec:EConvex}
The goal of this section is to build the fundamental tools necessary to adapt basic convex theory to our equivariant scenario.
\begin{defin}
	An action of a group $G$ on a manifold $X$ is said to be \emph{free} if $g\cdot x = x$, for some $g \in G$ and $x \in X$, implies $g = \id_X$.
\end{defin}

\begin{prop}\label{Gtrasv}
	Fix a finite group $G$; let $M$, $N$ be a $G$-manifolds, the latter equipped with a $G$-invariant Riemannian metric, $A$ a $G$-invariant closed subset of $M$, and $Y$ a $G$-submanifold of $N$. Let $f\colon M \to N$ be a smooth equivariant map transversal on $A$ to $Y$ in $N$. Suppose the $G$-action on $M \setminus A$ is free. Then for an arbitrary $G$-invariant positive continuous function $\delta \colon M \to \R$ , there exists a smooth $G$-map $g \colon M \to N$ satisfying the following conditions:
	\begin{enumerate}
	\item $g$ is transversal on $M$ to $Y$ in $N$;
	\item $g|_{A}=f|_{A}$;
	\item $d_{N}(f(x),g(x))<\delta(x)$ for all $x \in M$, where $d_{N}$ stands for the distance function on $N$ induced from the Riemannian metric of $N$.
	\end{enumerate}
\end{prop}
\begin{proof}
	See \cite[Thm.\ 1.1]{morimotoEQUIVARIANTTRANSVERSALITYTHEOREM2014}.
\end{proof}
\begin{lemma}\label{nondegequiv}
	Let $S$ be a closed surface, equipped with a involution that is free outside a finite number of points, and with an anti-invariant vector field with no closed orbits, that does not vanish on such points. Then it can be $C^\infty$-approximated by an anti-invariant vector field having a finite number of singularities, all non-degenerate (i.e.\ transverse to the zero-section).
\end{lemma}
\begin{proof}
	Set $G \coloneqq \Z/2\Z$. Let's consider the $G$-spaces $M \coloneqq S$ and $N \coloneqq TS$, with the action on the second space being
	\[
		g\cdot(p,v) = (g\cdot p, - d_pg[v])\,.
	\]
	The anti-invariant vector field are precisely the $G$-equivariant sections in the sense of \Cref{Gtrasv}. Moreover, we know that on the fixed points of $G$ the field $X$ does not vanish. Hence we have that
	\begin{enumerate}
		\item the $G$ action outside such points is free;
		\item on such points $X$ is vacuously transversal to the zero-locus of $TS$.
	\end{enumerate}
	
	Hence we can approximate $X$ with a new map $X' \colon S \to TS$ that is transversal to the zero-section, but that is not necessarily a section. Let's set $X'(x) = (h(x), Y(x))$ to be the vector $Y(x)$ belonging to the tangent space of $h(x)$. Set $\pi$ to be the natural projection from $TS$ to $S$.
	
	Since $h$ it is arbitrarily close to the identity map, it can be supposed to be a diffeomorphism of $S$ to itself. Then $Z \coloneqq X' \circ h^{-1}$ is vector field close to $X$ and transverse to $S$. Moreover
	\[
		(h \circ u)(x) = (\pi \circ X' \circ u)(x) = \pi[- d_{h(x)}u(Y(x))] =  u(h(x))\,,
	\]
	hence $h$ is $u$-equivariant. Finally, $Z$ is an anti-invariant vector field:
	\[
	\begin{split}
		(Z \circ u)(x) &= (X' \circ h^{-1} \circ u)(x)\\
		&= (X' \circ u \circ h^{-1})(x) \\
		&=- d_{hh^{-1}(x)}u[X'(h^{-1}(x))]\\
		&= - d_xu[Z(x)]\,.
	\end{split}
	\]
	So $Z$ is the required perturbation.
\end{proof}

Recall that given a vector field $X$ on a surface $S$, a singularity $p$ is non-degenerate iff, red in charts, the linear map $d_pX$ from $\R^2$ to itself non-zero determinant. Such a point is necessarily an isolated singularity. Moreover, if such linear map has eigenvalues with non-zero real part, then $p$ is \emph{generic}. Finally, if the real parts of such eigenvalues have the same signs, we say that $p$ is \emph{elliptic}. If not, $p$ is \emph{hyperbolic}.
\begin{defin}\label{MS}
	A vector field $X$ is said to be of Morse--Smale type if the following occur:
	\begin{enumerate}
		\item there is only a finite number of singularities, all generic;
		\item the $\alpha$- and $\omega$- limit sets of every trajectory can only be singularities or closed orbits;
		\item no trajectory connects hyperbolic points;
		\item there is only a finite number of closed orbits, all ``simple''.
	\end{enumerate}
	We are not going to specify the meaning of the last requirement, since we will avoid close orbit all together.
\end{defin}

The next proposition is an equivariant adaptation of classical Peixoto's work regarding structurally stable vector fields on surfaces \cite{peixotoStructuralStabilityTwodimensional1962}.
\begin{prop}\label{CinftyMorseSmale}
	Let $(S, \Omega)$ be an oriented closed surface equipped with an area form. Moreover, suppose that $S$ is equipped with a $\Z/2\Z$-action having a finite number of fixed point and that preserves $\Omega$, and with an anti-invariant vector field with no closed orbits, that does not vanish on such points. Then it can be $C^\infty$-approximated by a anti-invariant vector field that is Morse--Smale.
\end{prop}
\begin{proof}
	Thanks to \Cref{nondegequiv} we can perturb $X$ such that all the singularity are non-degenerate. Then, an adequate perturbation allows us to make such singularities generic; since such modification is done in neighborhoods of the singularities, which are not fixed points, we can simply work symmetrically in order to obtain an anti-invariant perturbation.
	
	Regarding the second and the third requests of \Cref{MS}, Peixoto achieves them by perturbing $X$ away from singularities, in arbitrarily small rectangles. Those are rectangles defined around fixed point, where the $x$-coordinate is parameterized using the flow of $X$. Hence, using the same procedure as of \Cref{tubsurf}, we can fix coordinates $(x,y)$ where the $u$-action is given by $(-x,-y)$.
	
	Thus, since all Peixoto's perturbation are of the form
	\[
		X' = X + \epsilon\, u\, \phi\, Z
	\]
	with 
	\[
	\begin{cases}
		\epsilon, u \in \R\\
		\phi \text{ some map } S \to \R, \phi > 0 \text{ only on the rectangle}\\
		Z = (0,\pm1)
	\end{cases}\,,
	\]
	it is immediate that $X'$ is anti-invariant.
	
	Finally, close orbits cannot appear by small perturbations.
\end{proof}

We want to use \Cref{CinftyMorseSmale} to perturb spheres such that the characteristic foliation become of Morse--Smale type. Let's start assuring that we have the right hypothesis.
\begin{lemma}\label{LefschetzS2}
	Let $S$ be a closed surface isomorphic to $S^2$ equipped with an involution $g$ such that, at every fixed point, the differential is minus the identity. Then $g$ must have exactly two fixed points.
\end{lemma}
\begin{proof}
	Consider the so-called \emph{Lefschetz number}
	\[
		\Lambda_g \coloneqq \Tr(g_{*,0}) + \Tr(g_{*,2})\,,
	\]
	where $g_{*,i}$ are the map induced on $H_i(S^2,\mathbb{Q})$, which is isomorphic to $\mathbb{Q}$ for $i = 0,2$ and vanishing otherwise.
	
	The trace $\Tr(g_{*,0})$ is $1$ since $S$ is connected. The trace $\Tr(g_{*,1})$ is $1$ since $g$ is orientation-preserving diffeomorphism. Hence $\Lambda_g$ is $2$.
	
	By using classical facts about fixed point theory, that can be found in \cite[Ch. VII]{albrechtlecturesonalgebraictopology}, one learns that
	\begin{enumerate}
		\item all the fixed point have \emph{fixed point index} $1$, since around such points $u$ can be viewed as minus the identity,
		\item since such fixed points are isolated, they are finite and the Lefschetz number $\Lambda_g$ is the sum of their index.
	\end{enumerate}
	Hence $g$ admits exactly 2 fixed points.
\end{proof}
\begin{prop}\label{CinftysurfMorseSmale}
	Given an $u$-oriented surface $S \subseteq \R^3$ diffeomorphic to $S^2$, and $\xi$ a tight invariant contact structure in $\R^3$, we can equivariantly $C^\infty$-perturb it such that it's characteristic foliation becomes Morse--Smale. 
\end{prop}
\begin{proof}
	Let's firstly observe that the condition
	\[
		\diver_{\Omega}(X)(p) \neq 0 \text{ if } X(p) = 0
	\]
	depends only on the field $X$, and is open in the $C^\infty$ topology. Moreover, since $\xi$ is tight, and $S$ is diffeomorphic to $S^2$, the characteristic foliation cannot have closed orbits, since the interior would be a disk violating the tightness condition \cite[Prop.\ 4.6.28]{geigesIntroductionContactTopology2008}.
	
	Due to \Cref{LefschetzS2}, the involution $u|_S$ has exactly 2 fixed points. Moreover, at such of a point $p$ the planes $\xi_p$ and $T_pS$ cannot coincide, since $d_pu$ reverses the orientation of the former, while keeping the one of the latter. Hence we are in the setting for applying \Cref{CinftyMorseSmale}. There exists another vector field $X' \in \XX(S)$, that is Morse--Smale, anti-invariant, and $C^\infty$-close to $X$. By \Cref{charinducescontact} we can find an anti-invariant contact form $\alpha'$, on a close neighborhood of $S$, inducing $X'$. Moreover it is $C^\infty$-closed to $\alpha$, and it can be made to coincide with the latter away from $S$.
		
	Hence we can define $\alpha_t$ as $t\alpha' + (1-t)\alpha$, which is a contact form since $\alpha$ and $\alpha'$ are sufficiently closed. Thanks to \Cref{EGray} we can find an equivariant isotopy $\psi_t$ of $\R^3$ such that $T\psi_t(\xi_t) = \xi$ \footnote{To be precise, \Cref{EGray} produces $\psi_t^{-1}$.}. Observe that $\psi_1^{-1}(S)$ has the desired characteristic foliation. In fact, if we denote with $\phi$ the restriction of $\psi_1$ on $\psi_1^{-1}(S)$, and if we impose the area form $\phi^*\Omega$ on $\psi_1^{-1}(S)$, then 
	\[
		i_{\phi^*X'}\, \phi^*\Omega = \phi^*(i_{X'}\Omega) = \phi^*j^*\alpha_1 = k^*\psi_1^*\alpha_1 = f\,k^*\alpha\,,
	\]
	with $j$ and $k$ the inclusion of $S$ and $\psi^{-1}_1(S)$ in $\R^3$, and $f$ a positive function.
\end{proof}

Let's enter in the main realm of this section: the one of convex surfaces. Recall that a \emph{contact vector field} with respect to a contact structure $\xi$ is a field such that its flow preserves the plane distribution. Moreover the set of smooth functions from $\R^3$ (or $S^3$) to $\R$ and the one of contact vector field are in bijection through the following assignments \cite[Thm.\ 2.3.1]{geigesIntroductionContactTopology2008}:
\[
\begin{split}
	 X &\mapsto \alpha(X)\\
	 H &\mapsto X_H \text{ such that } \alpha(X_H) = H \text{ and } i_{X_H} d\alpha = dH(R_\alpha)\alpha-dH.	 
\end{split}
\]
\begin{defin}
	A $u$-oriented compact surface $S$ is said to be \emph{convex} if there exists a contact vector field which is invariant, defined near and transverse to $S$.
\end{defin}
\begin{prop}
	An $u$-oriented surface $S$ is convex with respect to an invariant contact structure $\xi$ if and only if there exists an equivariant embedding $S \times \R \to \R^3 $ such that $\Psi^*\alpha$ determines a vertically invariant contact structure which is $u$-invariant.
\end{prop}
\begin{proof}
	Let's consider the anti-invariant function $H \coloneqq \alpha(X)$, with $X$ the contact transverse vector field. Let's perturb $H$ such that it vanishes away from $S$ by remaining anti-invariant. The new contact vector field $Y$ associated to $H$ is invariant, and
	\[
		\Psi(x,t) \coloneqq \psi_t^Y(p)
	\]
	is the required embedding. Moreover, since $u^*Y = Y$, then $\psi_t$ is $u$-equivariant in the spacial argument, i.e.\ $\Psi$ is equivariant in the sense of our convention on $S \times \R$. Hence $\Psi^*\alpha$ is anti-invariant. And since
	\[
		T\Psi(p,t)(\partial_t)= \dot{\psi}^Y_t(p)= Y(\psi^Y_t(p))\,,
	\]
	we have that $\Ker(\Psi^*\alpha)$ is vertically invariant, since $Y$ is a contact vector field.
	
	For the converse, just set $Y \coloneqq T\Psi(\partial_t)$.
\end{proof}
\begin{prop}
	In the previous situation we can suppose that $\Psi^*\alpha$ is a vertically invariant, and anti-invariant, contact form.
\end{prop}
\begin{proof}
	Keeping the notation of the previous proof, we just proved that $\Lie_{\partial_t}(\Psi^*\alpha) = \mu \Psi^*\alpha$, with $\mu\circ u$ coinciding with $\mu$. And since
	\[
		\Lie_{\partial_t}(\lambda \Psi^*\alpha) = \biggl(\frac{\dot\lambda}{\lambda} + \mu\biggr)\cdot(\lambda\Psi^*\alpha)
	\]
	holds, we just need to adjust the form by picking 
	\[
		\lambda \coloneqq \exp\Bigl(-\int_0^t\mu(p,s)\,ds\Bigr)\,.
	\]
	
	Moreover
	\[
		\lambda \circ u = \exp\Bigl(-\int_0^t(\mu \circ u)(p,s)\,ds\Bigr) = \lambda\,.
	\]
	
	Hence $\lambda\Phi^*\alpha$ is a contact form with the desired properties.
\end{proof}
As in the non-equivariant scenario, the information about a convex surface is contained in ``small'' family of circles, called \emph{dividing curves}.
\begin{defin}\label{defconvex}
	Let $\FF$ be a singular foliation on a $u$-oriented closed surface $S$. An invariant collection $\Gamma$ of embedded circles is said to divide $\FF$ if
	\begin{enumerate}
		\item $\Gamma$ is transverse to $\FF$; in particular, it does not meet the singularities of $\FF$,
		\item there is an invariant area form $\Omega$ on $S$ and an anti-invariant vector field $X$ defining $\FF$ such that
		\[
		\Lie_X\Omega \neq 0
		\]
		on $S \setminus \Gamma$, and with $S_\pm = \{p \in S \mid \pm \diver_\Omega(X) > 0\}$. so that if $S \setminus \Gamma = S_+ \sqcup S_-$, the vector field $X$ points out of $S_+$ along $\Gamma$.
	\end{enumerate}
\end{defin}
\begin{prop}\label{divconvex}
	If a closed, $u$-oriented compact surface $S$ is divided by a collection $\Gamma$ of embedded circles, then $S$ is convex with respect to $\xi$.
\end{prop}
\begin{proof}
	Assume that $S_\xi$ is divided by some set $\Gamma$. We need to show that $S$ is convex. Let $\Omega$ and $X$ be the area form and vector field, respectively, from \Cref{defconvex}.
	
	Set $\beta = i_X\Omega$ and consider the $\R$-invariant 1-form $\alpha = \beta + w\,dz$ on a neighborhood $S \times \R$ of $S \times \{0\}$ in $M$, with $w$ an arbitrary anti-invariant smooth function on $S$, and with $z$ denoting the $\R$-coordinate. The singular foliation induced by $\alpha$ on $S$ is precisely $S_\xi$. The condition for $\alpha$ to be a contact form is given by Equation \eqref{eq:contcond}, that in this context becomes
	\[
	w \diver_\Omega(X) - X(w) > 0.
	\]
	On $S_\pm$, this condition is satisfied for $w \equiv \pm 1$.
	
	Now use the flow of $X$ to identify a neighborhood $A$ of $\Gamma$ in $S$, made up of a disjoint collection of annuli, with $\Gamma \times [-\epsilon, \epsilon]$, such that $\Gamma$ becomes identified with $\Gamma \times \{0\}$.
	
	Using again Bredon's result about invariant tubular neighborhoods, we can suppose that the previous diffeomorphism is such that every $u$-invariant component of $\Gamma$ is identified with a neighborhood of $S^1$, with $u$-action given by
	\[
		(\theta, s) \mapsto (-\theta,-s)\,.
	\]
	
	The parameter $s \in [-\epsilon, \epsilon]$ is chosen such that $X = -\partial_s$, so that the sign of $s$ accords with the sign in $S_\pm$. Now set
	\[
	h(\theta, s) \coloneqq \exp\left(-\int_0^{s} \diver_\Omega(X)(\theta, t)\,dt\right)
	\]
	and make the Ansatz
	\[
	w(\theta, s) = g(\theta, s) \cdot h(\theta, s).
	\]
	Then, observing that $X(w) = -\partial w/\partial s$, one can compute that
	\[
	w \diver_\Omega(X) - X(w) = \frac{\partial g}{\partial s}\,h(\theta, s).
	\]
	Therefore, $\alpha$ will be a contact form, coinciding away from $\Gamma$ with $\beta + dz$, if we choose $g$ to satisfy
	\begin{equation}\label{casesg}
	\begin{cases}
		\partial g/\partial s > 0, \\
		g(\theta, s) = \pm 1/h(\theta, s) \quad \text{near } s = \pm\epsilon.
	\end{cases}
	\end{equation}
	
	Such a smooth function $g(\theta,s)$ can indeed be found, at least if we ignore the $u$-action. We have thus constructed an $\R$–invariant contact structure $\Ker\alpha$ on $S\times \R$,	so $S$ is a convex surface with respect to this new contact structure. Moreover, since $\Ker\alpha$ induces the given characteristic foliation $S_\xi$, the germs of $\Ker\alpha$ and $\xi$ near $S$ coincide by \Cref{EGiroux}. This means that $S$ is also a convex surface for $\xi$.
	
	Moreover, on the invariant component of $\Gamma$ we have that $w \circ u$ coincides with $-w$. Indeed,
	\[
	\begin{split}
		h \circ u &= \exp\left(-\int_0^{-s} \diver_\Omega(X)(-\theta, t)\,dt\right) \\
		&= \exp\left(-\int_0^{s} -\diver_{\Omega}(X)(-\theta, -t)\,dt\right)\\
		&= \exp\left(-\int_0^{s} - u^*\diver_\Omega(X)(\theta, t)\,dt\right)\\
		&= \exp\left(-\int_0^{s} - \diver_{u^*\Omega}(u^*X)(\theta, t)\,dt\right)\\
		&= \exp\left(-\int_0^{s} \diver_\Omega(X)(\theta, t)\,dt\right) = h\,,
	\end{split}
	\]
	and by choosing $g$ to be odd with respect to $s$, we have that $g$, and hence $w$, is anti-invariant. Regarding the components exchanged by $u$, one just need to define the function $w$ one of the coupled component, and set $w \circ u$ as $-w$. 
	
	Finally, we can additionally  suppose that $g(p,0) = 0$, so that $\Gamma$ is actually the \emph{dividing set} with respect of the vector field $\partial_z$; see \cite[Def.\ 4.8.1]{geigesIntroductionContactTopology2008} for the relative definition.
\end{proof}
In the case of $S$ being a sphere, the tightness condition forces the dividing set to be as simple as possible.
\begin{prop}\label{divsetS2iscircle}
	Let $S$ be a closed, convex $u$-oriented surface. If a dividing set $\Gamma$ contains a circle $\gamma$ contractible in $S$, then any vertically invariant neighborhood of $S$ is overtwisted, unless $S = S^2$ and $\Gamma$ is connected.
\end{prop}
\begin{proof}
	See \cite[Prop.\ 4.8.13]{geigesIntroductionContactTopology2008}.
\end{proof}

With \Cref{divconvex} in mind, we want to prove that in the case of $S$ being a sphere and contact structure being tight the convex condition for a surface is in fact generic. Before doing that, we need some definition regarding surfaces isomorphic to annuli.
\begin{defin}
	Let $A$ be an annulus $S^1 \times [-1,1]$ equipped with a orientation preserving $\Z/2\Z$-action $u$. A \emph{nerve} $N$ is an invariant set of the form $\{(\theta, f(\theta))\}$, where $f$ is a smooth function $S^1 \to [-1,1]$, such that the action on $N$ is given by $\theta \mapsto -\theta$.
\end{defin}
\begin{lemma}\label{nerve}
	Suppose that the action $u$ satisfies the following properties:
	\begin{enumerate}
		\item at every fixed point the differential, viewed as an endomorphism of the tangent space, is minus the identity;
		\item the action can be written as $(\theta,t) \mapsto (g(\theta), h(\theta,t))$;
		\item $h(\theta,\cdot)$ reverses the orientation of the interval for every $\theta$.
	\end{enumerate}
	Then $A$ has a nerve.
\end{lemma}
\begin{proof}
	Let's use the same ideas of \Cref{LefschetzS2}.
	
	Since $u$ is orientation-preserving, and since $h$ is orientation-reversing, then the map $g \colon S^1 \to S^1$ is an orientation-reversing diffeomorphism. Moreover, the Lefschetz number
	\[
		\Lambda_u \coloneqq \Tr(u_{*,0})- \Tr(u_{*,1})\,
	\]
	coincides with $\Lambda_g$. In fact $A$ is homotopically equivalent to $S^1$, with $u$ getting identified with $g$. 
	
	The trace $\Tr(g_{*,0})$ is $1$ since $A$ is connected. The trace $\Tr(g_{*,1})$ is $-1$ since $g$ is homotopy equivalent to a reflection, being orientation-reversing. Hence $\Lambda_u = \Lambda_g$ is $2$.
	
	Thus $u$, and $g$, admits exactly 2 fixed points. Let $(\theta_1, t_1)$ and $(\theta_2, t_2)$ be such fixed points. Up to rescaling, one can consider them to be $(0,t_1)$ and $(\pi,t_2)$. Finally, if we consider an arbitrary map $m \colon [0,\pi] \to [-1,1]$ a nerve is given by the set
	\[
		N \coloneqq \Imm(m) \cup u(\Imm(m))\,. 
	\]
	
	Since the differential action is minus the identity, the set $N$ glues \emph{a fortiori} to a smooth submanifold.
\end{proof}
\begin{prop}\label{MorseSmaleconvex}
	A closed surface $S$ that is $u$-oriented, with Morse--Smale foliation, is convex.
\end{prop}
\begin{proof}
	By \Cref{divconvex} it suffices to show that $S_\xi$ admits dividing curves. We retain the vector field $X$ defining $S_\xi$, but we are going to adjust the area form $\Omega$ on $S$.
	
	Let $S^0_\pm \subset S$ be submanifolds with boundary made up of discs around the elliptic points (of the corresponding sign) of $S_\xi$, annuli around the repelling or attracting periodic orbits, respectively, and bands around the stable or unstable separatrices of the positive or negative hyperbolic points, respectively. We may assume that $S^0_+ \cap S^0_- = \emptyset$ and that $X$ is transverse to the boundary of $S^0_\pm$.
	
	On $A := S \setminus (S^0_+ \cup S^0_-)$ the characteristic foliation is non-singular, transverse to the boundary, and without closed leaves. Hence $A$ is a collection of annuli; let's fix a diffeomorphism 
	\[
		h \colon A \to \bigsqcup_i S^1 \times [-1,1] \eqqcolon \bigsqcup_i S_i\,.
	\]	
	For every annulus $S_{\bimath}$ preserved by $u$ \Cref{nerve} fixes a nerve, and sets a diffeomorphism such that the action on $S_{\bimath}$ is of the form $(\theta,s) \mapsto (-\theta, -s)$.
	
	Moreover, we can suppose the sets $S_\pm^0$ to be $u$-invariant, by choosing $S^0_-$ and setting $S^0_-$ as $u(S^0_+)$. In fact, $u$ must reverse the sign of the singular points, since it keeps the orientation of the surface $S$ while reversing the one of the contact planes.
	
	On $A$ we impose an area form $\Omega_A$ of the form $f(s)\,d\theta\wedge ds$ such that $\pm\diver_{\Omega_A}(X) > 0$ on $A\cap S_\pm$. Indeed, one computes 
	\[
	\diver_{\Omega_A}(-\partial_s)=-\frac{f'}{f}\,,
	\]
	so we simply have to choose a positive function $f(s)$ that is strictly increasing	on $[-\epsilon,0)$ and strictly decreasing on $(0,\epsilon]$. Moreover we can evidently choose such $f$ such that is is even. In this way $u^*\Omega_A$ coincides with $\Omega_A$.
		
	Equation
	\[
		\diver_{g\Omega}(X)=\frac{X(g)}{g} +\diver_\Omega(X)
	\] 
	shows that the divergence at the singular points will not change if we rescale the area form to $g\Omega$, with $g: S \to \R^+$. The same equation shows that we may achieve $\operatorname{div}_{g\Omega}(X) > 0$ on $S^0_+$ by having $g$ grow fast enough along the flow lines of $X$ (and equal to some constant near each singular point). Likewise, we can ensure that $\operatorname{div}_{g\Omega}(X) < 0$ on $S^0_-$. If we suppose $f$ being $s$-even, we can take $g$ to be invariant.
	

	Let's now use the $s$ coordinate to patch $\Omega_A$ with $g\Omega$ to define a third volume form $\Omega'$. Evidently we can make the patching $u$-invariant by using a an even bump function.
	
	
	 Finally let's define $\Omega'' \coloneqq k\Omega'$, with $k$ a function that is $\equiv 1$ at near the boundary of our annuli, $u$-invariant, and equal to some large positive constant around $\Gamma$. If we make it grow fast enough, then $\pm\diver_{k\Omega''}(X) > 0$ on the complementary of $\Gamma \times [-\epsilon/3, \epsilon/3]$. Finally if the divergence condition of \Cref{defconvex} is achieved with $\Omega_A$, than it must be achieved also with $K\Omega_A$, hence also for $\Omega''$ near $\Gamma$.
	
	
	In conclusion an arbitrary nerve of $A$ divides $S_\xi$, and $S$ is convex by \Cref{divconvex}.
\end{proof}
\begin{cor}\label{Cinftyconvex}
	Given an $u$-oriented surface $S \subseteq \R^3$ diffeomorphic to $S^2$, and $\xi$ a tight invariant contact structure in $\R^3$, we can equivariantly $C^\infty$-perturb $S$ such that it becomes convex.
\end{cor}
\begin{proof}
	Just use \Cref{CinftysurfMorseSmale} to make $S_\xi$ of type Morse--Smale, so that $S$ is convex by \Cref{MorseSmaleconvex}.
\end{proof}
\section{Equivariant Tomography Theory}\label{sec:tomography}

In this section we will look at the \emph{tomography theory}, i.e.\ the study of contact structures on $S^2 \times [-1,1]$ via the characteristic foliation on the slices $S^z \coloneqq S^2 \times \{z\}$.

We suppose the product $S^2 \times [-1,1]$ to be equipped with an involution $u$ such that
\begin{enumerate}
	\item the fixed-point set is non-trivial, made of two segments each connecting $S^2 \times \{-1\}$ with $S^2 \times \{1\}$;
	\item at each fixed point the differential is diagonalizable as $[1] \oplus [-1]\oplus[-1]$.
\end{enumerate}
We will call such scenario to be an \emph{equivariant tomography}. Observe that in such scenario the orientation of the boundary spheres must be preserved, with the differential of the involution at a fixed point restricting to $TS^{\pm1}$ as minus the identity.

With such setting in mind, the aim of this section is to prove the following result.
\begin{thm}\label{thm:mainthmtomographyStwo}
	Let $S^2 \times [-1,1]$ an equivariant tomography. The set of invariant tight contact structure on $S^2 \times [-1,1]$ having fixed characteristic foliation on the boundary, up equivariant isotopy rel boundary, is a $\Z$-torsor.
\end{thm}

Before focusing to the tools required for proving \Cref{thm:mainthmtomographyStwo}, let's prove an important reduction result for our action. It will be preceded by a couple of technical lemmas.
\begin{lemma}\label{Palais}
	Given a compact manifold $M$, and a submanifold $N$, the natural map from $\Diff(M)$ to $\Emb(N,M)$ is a Serre fibration (where the spaces are equipped with the $C^\infty$-topology).
\end{lemma}
\begin{proof}
	Actually more holds: by Palais the above map is a locally trivial fiber bundle \cite{PalaisLocaltrivialityoftherestrictionmap}. However, it is useful, for future reference, to have a partial proof.
	
	Consider an arbitrary map $\phi_{u,t}$ from $D^n \times I$ to $\Emb(N,M)$, such that $\phi_{u,t}$ is a smooth isotopy of embeddings from $N$ to $M$; fix a lift $\tau_u$ of $\phi_0$. We want to extend such lift to the whole interval $I$. Said otherwise, we need to find an isotopy $\sigma_{u,t}$ of diffeomorphisms such that $\sigma_{u,t}(N)$ coincides with $\phi_{u,t}(N)$ and $\sigma_{u,0}$ coincides with $\tau_u$.¨
	
	For every point $u \in D^n$  we define the following vector field:
	\[
		X_{u,t} \coloneqq \frac{d}{dt} \phi_{u,t} \in TM|_N\,.
	\]
	By fixing a Riemannian metric on $M$, one can canonically extend $X_{u,t}$ to a new vector field $\tilde{X}_{u,t}$ on the whole manifold $M$, that vanishes outside a small neighborhood of $N$. Hence, we just have to integrate such vector field to define a new isotopy $\Phi_{u,t}$ of $M$ such that
	\[
	\begin{cases}
		\Phi_{u,0} = \id_M \\
		\frac{d}{dt} \Phi_{u,t} = \tilde{X}_{u,t}
	\end{cases}\,.
	\]
	
	Since $\tilde{X}_{u,t}$ extended $X_{u,t}$ the equality
	\[
		\Phi_{u,t} \circ \phi_{u,0} = \phi_{u,t}
	\]
	holds by construction. Hence we define $\sigma_{u,t}$ as $\Phi_{\bullet,\bullet} \circ \tau_{\bullet}$. Then $\sigma_{u,0}$ coincides with $\sigma_u$, and by the previous equality $\sigma_{u,t}$ lifts $\phi_{u,t}$.
	
	Palais' argument is similar, and quite easy to understand. He uses vector fields and exponential maps to find local section of the map he claims to be locally trivial. Observe that the above reasoning is not complete, since \emph{a priori} one has to deal with continuous maps from $D^n \times I$ to $\Emb(N,M)$, without any hypothesis on smoothness in times. However, if $N$ and $M$ are low-dimensional, I think one could be able to slightly modify the above proof to be correct.
\end{proof}
\begin{prop}\label{InvolutionsS2}
	Consider the set $\mathcal{I}$ of involution of $S^2$ fixing two fixed points $s$, $s'$. Then the natural conjugation action, from $\Diff(S^2, s, s')$ to $\mathcal{I}$, is a Serre fibration.
\end{prop}
\begin{proof}
	The conjugation action is surjective, see \cite[Thm.\ 1]{EilenbergSurlestransformationsperiodiques}. However, more is true, that can be seen by delving into the cited proof. Let's fix an arch $A$ from $s$ to $s'$, and let's consider an involution $g \in \mathcal{I}$. Then the diffeomorphism conjugating $g$ to a reference involution $g_0$ can be easily build. Firstly, let's fix a standard map bringing $A \cup g(A)$ to $A \cup g_0(A)$. Then, the arbitrary comes into extending such map to diffeomorphism between the disks bounded by $A \cup g(A)$ and $A \cup g_0(A)$. At this point is simply a matter of extending embeddings, and we can conclude with Palais' work.
	
	Observe that if we start with a smooth isotopy $g_t \in \mathcal{I}$, then the proof of \Cref{Palais} tells us that we can find a smooth isotopy of diffeomorphisms $\psi_t$ fixing $s$ and $s'$ such that $\psi_t \circ g_t \circ \psi_t^{-1}$ coincides with $g_0$.
\end{proof}
\begin{prop}\label{standardizedtomography}
	There exists a diffeomorphism from $S^2 \times [-1,1]$ to itself that conjugates $u$ to an involution mapping $(s,z)$ to $(u|_{S}(s), z)$, where $S = S^2 \times \{-1\}$.
\end{prop}
\begin{proof}
	The properties of $u$ tell us that the fixed point set is a $1$-dimensional submanifold, and that the quotient $Q$ of $M \coloneqq S^2 \times [-1,1]$ by $u$ is a smooth orientable 3-manifold, with boundary given by the disjoint union of the two spheres coming from $M$.
	
	Moreover, Armstrong's Theorem for orbifolds \cite{ArmstrongThefundamentalgroupoftheorbitspaceofadiscontinuousgroup} can be applied: the quotient has fundamental group isomorphic to $G/H$, where $G = \Z/2\Z$ and $H$ is the subgroup made by the elements having at least one fixed point. Hence such fundamental group is trivial.
	
	So $Q$ is a simply connected 3-manifold with boundary $S^2 \sqcup S^2$. Thanks to the $h$-cobordism theorem between surfaces, holding due to Poincaré Conjecture's being proven, the manifold $Q$ is actually isomorphic to $S^2 \times [-1,1]$. Let $\phi$ be a diffeomorphism from $Q$ to such product manifold.
	
	Let's set $h$ to be the composition $p_2 \circ\phi \circ \pi$, where $\pi \colon M \to Q$ is the product map and $p_2$ is the projection onto the second factor.
	
	Such map is a surjective submersion, and $M$ is compact. Hence due to Ehresmann's Lemma \cite{EhresmannLesconnexionsininitesimales} such map is a locally trivial fiber bundle over $[0,1]$. Moreover, since such space is contractible, there exists a diffeomorphism of fiber bundles $\psi$ from $M$ to the trivial fiber bundle $F \times [0,1]$, where $F = S^2$.
	
	I claim that $\psi$ is the diffeomorphism we are looking for. In fact, consider a point $(s,z)$ in $S^2 \times [-1,1]$, and apply $\psi \circ u \circ \psi^{-1}$. The points $\psi^{-1}(s,z)$ and $u(\psi(s,z))$ are in the same $h$-fiber. Since $\psi$ sends $h$-fibers in $M$ to sets of the form $S^2 \times \{z\}$, then $\psi(u(\psi^{-1}(s,z)))$ must have the $z$-coordinate unchanged.
	
	So we arrived at a situation in which $u(s,z)$ is of the form $(g(s,z), z)$. We want to find a diffeomorphism of $M$ making $g$ independent of $z$. Firstly, since $u|_{S^z}$ preserves the orientation at the $z = \pm1$ levels, by continuity it must do as at every level. Hence, the fixed point set must now be made of two segments intersecting each level exactly ones: if not, a sphere $S^2 \times \{z\}$ would be equipped with an orientation-preserving involution having only regular fixed points, and such fixed point being more than two. This contradicts the computation done in \Cref{LefschetzS2}.
	
	As a consequence, we can consider the smoothly varying pair $(s_z,s'_z)$ of such fixed point. By the proof of \Cref{Palais} there exists a smooth family of diffeomorphisms $\psi_z$ such that $\psi_z(s_z, s'_z)$ equals $(s_0,s'_0)$. Hence, the diffeomorphisms $\psi_z \times \id_{[-1,1]}$ maps $M$ to itself, and $u$ to an involution having segments $\{s\} \times [-1,1]$ and $\{s'\} \times [-1,1]$ as fixed points.
	
	Finally, due to \Cref{InvolutionsS2} we can fix fix a smooth family of diffeomorphisms $\psi_z$ such that $\psi_z \circ g_z \circ \psi_z^{-1}$ coincides with $g_0$. The final diffeomorphism from $M$ to itself is then given by $\psi_z \times \id_{[-1,1]}$.
\end{proof}

So, up to changing coordinates, we can suppose that $u$ is of the form prescribed by \Cref{standardizedtomography}.
\subsection{Reconstruction and Uniqueness Lemmas}
Before dwelling into the proof of \Cref{thm:mainthmtomographyStwo}, let's state a couple of lemmas, that will indicate the direction we are trying to follow.
\begin{lemma}[Reconstruction Lemma]
	Let $S \times [-1,1]$ be an equivariant tomography, equipped with two invariant contact structures $\xi_0$ and $\xi_1$, such that the foliations $S^z_{\xi_i}$ coincide for each $z \in [-1,1]$. Then $\xi_0$ and $\xi_1$ are equivariantly isotopic rel boundary.
\end{lemma}
\begin{proof}
	The assumption on the characteristic foliations implies that we can write
	\[
		\alpha_i = \beta_z + w^i_z\,dz\,,
	\]
	with $\beta_z$ independent of $i = 0,1$. Moreover, the condition
	\[
		 w_z^i d\beta_z +\beta_z \wedge (dw_z^i - \dot{\beta}_z) > 0
	\]
	is convex if $\beta_z$ is fixed. Hence
	\[
		\alpha_t \coloneqq t\alpha_1 + (1-t)\alpha_0
	\]
	is a contact form. It then follows from \Cref{EGray} that $\xi_0$ and $\xi_1$ are equivariantly isotopic; an isotopy rel boundary can be build as in \Cref{Girouxbound}, since the homotopy of contact forms is stationary on $TS^{\pm 1}$.
\end{proof}
\begin{lemma}[Uniqueness Lemma]\label{uniqueness}
	Let $S \times [-1,1]$ be an equivariant tomography, equipped with two invariant contact structures $\xi_0$ and $\xi_1$ with the following properties:
	\begin{enumerate}
	\item the characteristic foliations $S^{\pm 1}_{\xi_i}$ on the boundary of $S \times [-1,1]$ coincide for $i =0,1$;
	\item each surface $S^z$, $z \in [-1,1]$, is convex for either contact structure, and there is a smoothly varying family of multi-curves $\Gamma_z$ dividing both $S^z_{\xi_0}$	and $S^z_{\xi_1}$.
	\end{enumerate}
	Then $\xi_0$ and $\xi_1$ equivariantly are isotopic rel boundary
\end{lemma}
\begin{proof}
	Choose contact forms 
	\[
		\alpha^i = \beta^i_z + w^i_z\, dz, \quad i=0, 1,
	\]
	defining $\xi_i = \Ker \alpha^i$, with $\beta^0_{\pm 1} = \beta^1_{\pm 1}$. Recall the contact condition 
	\[
		A^i \coloneqq w^i_z\, d\beta^i_z + \beta^i_z \wedge (dw^i_z - \dot{\beta}^i_z) > 0, \quad i=0, 1,
	\]
	with the dot denoting the derivative with respect to $z$. The function $A^i$ is invariant.
	
	The convexity assumption allows us to choose, for each $z \in [-1, 1]$, functions $v^0_z$ and $v^1_z$ on $S$ that vanish exactly on $\Gamma_z$, only to first order, and that satisfy the contact condition
	\[
		B^i \coloneqq  v^i_z\, d\beta^i_z + \beta^i_z \wedge dv^i_z > 0, \quad i=0, 1.
	\]
	The function $B^i$ is invariant.
	
	The construction in the proof of Theorem 4.8.5 gives us two families of functions $v^i_z$ varying smoothly with $z$, and with $v^0_{\pm 1} = v^1_{\pm 1}$. Notice that for each $z$ the function $v^0_z/v^1_z$, defined a priori on $S \setminus \Gamma_z$ only, extends unambiguously to a positive function on all of $S$. The computation
	\[
	\begin{split} 
		v^0\, d\left(\frac{v^0}{v^1}\beta\right) + \left(\frac{v^0}{v^1}\beta\right) \wedge dv^0 &= \left(\frac{v^0}{v^1}\right)^2 v^1d\beta + \frac{v^0}{v^1} dv^0 \wedge \beta \\ 
		& \quad - \left(\frac{v^0}{v^1}\right)^2 dv^1 \wedge \beta + \frac{v^0}{v^1}\beta \wedge dv^0 \\ 
		&= \left(\frac{v^0}{v^1}\right)^2 (v^1 d\beta + \beta \wedge dv^1) 
	\end{split}
	\]
	shows that if we replace $\beta^1_z$ by $(v^0_z/v^1_z)\cdot \beta^1_z$ we may take $v^1_z= v^0_z \eqqcolon v_z$. Observe that after this replacement we still have $\beta^0_{\pm1} = \beta^1_{\pm1}$.
	
	With $\lambda \in \R^+$ set the anti-invariant form
	\[
		\alpha^i_t \coloneqq \beta^i_z + (1-t)w^i_z + t\,\lambda\, v_z\, dz, \quad i =0,1, \quad t\in [0,1]\,.
	\]
	Then
	\[
		\alpha^i_t \wedge d\alpha^i_t = ((1-t)A^i  - t\beta^i_z \wedge \dot{\beta}^i_z +t\lambda B^i) \wedge dz\,,
	\]
	which is positive for all $t \in [0,1]$ if $\lambda$ is chosen sufficiently large. Finally, set the family of anti-invariant forms $\alpha_s$ as
	\[
		\alpha_s \coloneqq (1-s)\alpha^0_1 +s \alpha^1_1 =(1-s)\beta^0_z +s\beta^1_z +\lambda\, v_z\, dz,\quad s \in  [0,1]\,.
	\]
	Then
	\[
		\alpha_s \wedge d\alpha_s = \lambda((1-s)B^0 +sB^1) \wedge dz+O_\lambda(1)\,,
	\]
	which is likewise positive for all $s \in [0,1]$, provided $\lambda$ is sufficiently large.
	Then
	\[
	\begin{cases}
	\alpha^0_{3t}, & 0 \le t \le 1/3, \\
	\alpha_{3t-1}, & 1/3 \le t \le 2/3,\\
	\alpha^1_{3-3t}, & 2/3 \le t \le 1,
	\end{cases}
	\]
	defines a homotopy via contact forms from $\alpha^0$ to $\alpha^1$ that is stationary on $TS^{\pm1}$. Now the proof concludes as in the reconstruction lemma.
\end{proof}
\begin{rmk}\label{rmkboundary}
	If in either of the two preceding Lemmas the contact structure $\xi_0$, $\xi_1$ coincide near the boundary, the proofs give isotopies that are stationary near the boundary.
\end{rmk}
Moreover, from the proof of the Reconstruction and Uniqueness Lemmas we can get the following result.
\begin{lemma}\label{parametricdividedetermined}
	Let $\xi$ be an invariant tight contact structure on an equivariant tomography $S^2 \times [-1,1]$, such that $S^z$ are $\xi$-convex. Then a smooth family of dividing circles $\Gamma_z$, coinciding near the boundary, is determined up to level-preserving isotopy $S^2 \times [-1,1]$ stationary near the boundary.
\end{lemma}
\begin{proof}
	Let $\Gamma_z^0$ and $\Gamma_z^1$ be two families of dividing circles.
	
	Similarly to the proof of the Uniqueness Lemma, we can write $\alpha$ as
	\[
		\alpha = 
		\begin{cases}
			\beta_z + v^0_z\,dz & \text{using the circle $\Gamma^0_z$}\\
			\beta_z + v^1_z\,dz & \text{using the circle $\Gamma^1_z$}
		\end{cases}\,.
	\]
	Hence we can define the family
	\[
		\alpha_t \coloneqq \beta_z + ((1-t)v^0_z + tv^1_z)\,dz\,,
	\]
	and appealing to the \Cref{Girouxbound} to obtain a level-preserving isotopy $\psi_t$ pulling back $\alpha_t$ to a multiple of $\alpha_0$. Since the zeros of $v^0_z$ and of any its positive multiple are the same, we get that $\psi_1$ maps the dividing curve $\Gamma^0_z$ to $\Gamma^1_z$.	
\end{proof}
\subsection{Defining the twisting torsor}
The aim of this subsection is to introduce define the correct twisting invariant appearing in the statement of \Cref{thm:mainthmtomographyStwo}. 





\begin{lemma}\label{twZtorsor}
	Let $F$ be a line in a 3-manifold $S \times [-1,1]$, intersecting each $S^z$ at a single point. Then the set of hyperplanes distribution $\xi_p \subseteq T_pM$ tangent to $F$, with fixed endpoints, up to isotopy rel boundary, is a $\Z$-torsor.
\end{lemma}
\begin{proof}
	The above equivalence classes are in bijection with the set of trivializations of normal bundle $NF \subseteq TM$ up to isotopy of $S \times [-1,1]$ rel boundary. It is then a standard fact in differential geometry that such object is a $\Z$-torsor.
\end{proof}

Let's now consider an arbitrary invariant tight contact structure $\xi$ on an equivariant tomography $S^2 \times [-1,1]$, and let's pick an arbitrary line $F = \{s\} \times [-1,1]$ of fixed points. For every point $p$ on $F$, the plane $\xi_p$ is preserved by $u$. Since $d_pu$ is diagonalizable as $[1] \oplus [-1] \oplus [-1]$, so must be $u|_{\xi_p}$. Finally, since $u$ reverses the orientation of the contact planes, $u|_{\xi_p}$ must necessarily be writable as $[1]\oplus [-1]$, and hence $F$ must be tangent to $\xi_p$.
\begin{defin}\label{twdif}
	Let $\xi_1$ and $\xi_2$ be two invariant tight contact structures on $S^2 \times [-1,1]$ coinciding at the boundary. Then, having chosen a line of fixed points $F$ , the \emph{$F$-twisting difference} is the $\Z$-difference of \Cref{twZtorsor} computed using $F$.
\end{defin}
The previous definition seems dependent on the line of fixed point we choose. A way to resolve that is to notice that in the case of a film of convex surfaces the quantity becomes independent. This is what we shall do.


Firstly, let's define the \emph{equivariant unparameterized embedded loop space} as
\[
	\Loop^{eq}(S^2) \coloneqq \Emb^{eq}(S^1,S^2)/\Diff^{eq,+}(S^1)\,,
\]
with the latter diffeomorphism group acting on the right. In the next \namecref{EBaer} will compute its fundamental group.
\begin{prop}\label{EBaer}
	The group $\Loop^{eq}(S^2)$ is weakly homotopically equivalent to $SO_2$, and its fundamental group is infinite cyclic.	
\end{prop}
\begin{proof}
	By quotenting $S^2$ by $u$ we obtain a double cover $S^2 \to S^2$, branched over two points $a$,  $b$. Hence the space we are looking to can be viewed as the space $E \coloneqq \Omega_{a,b}^{\emb}S^2$, i.e.\ the space of embedded paths from $a$ to $b$. 
	
	Let's fix an element $e_0 \in E$. Firstly we observe that there exists a map
	\[
	\begin{split}
		\phi \colon G \coloneqq \Diff^+(S^2, a,b) &\to \Omega^{\emb}_{a,b}S^2\\
		\phi &\mapsto \phi\circ e_0
	\end{split}
	\]
	with the former being the space of orientation-preserving diffeomorphism fixing the points $a$ and $b$. The surjectivity of $\phi$ comes from two results: since $\Omega_{a,b}S^2$ is connected, by Baer's theorem \cite{BaerIsotope} so must be the space $\Omega_{a,b}^{\emb}S^2$; then the isotopy extension theorem allows us to extend isotopies to paths of diffeomorphisms. Moreover, with the same approach of Palais, whose statement is given in \Cref{Palais}, one can show $\phi$ to be a Serre fibration.
	
	By cutting along $e_0$, we see that the subgroup of $G$ fixing $e_0$ can be identified with the space $\Diff^+(D^2, \partial D^2)$ of orientation-preserving diffeomorphism of the disk preserving the boundary. By Smale \cite{Smale1959} this space is contractible, and hence $E$ is weakly homotopically equivalent to the total space.
	
	Let's consider the following Serre fibration, by \Cref{Palais}:
	\[
		ev_a \colon \Diff^+(S^2, a) \to S^2 \setminus \{a\}
	\]
	given by evaluation at $b$. Since the base is contractible, the fiber $G$ is weakly homotopically equivalent to $\Diff^+(S^2,a)$. 
	
	Finally, by valuating at $a$ we get the final Serre fibration we are interested in:
	\[
		\Diff^+(S^2,a) \to \Diff^+(S^2) \to S^2\,.
	\]
	It sits in the following commutative diagram of fibrations:
	\[
	\begin{tikzcd}
		\Diff^+(S^2,a) & \Diff^+(S^2) & S^2 \\
		SO_2 & SO_3 & S^2
		\arrow[from = 1-1,to=1-2, hookrightarrow]
		\arrow[from = 1-2,to=1-3]
		\arrow[from = 2-1,to=2-2, hookrightarrow]
		\arrow[from = 2-2,to=2-3]
		\arrow[from = 2-1,to=1-1, hookrightarrow]
		\arrow[from = 2-2,to=1-2, hookrightarrow]
		\arrow[from = 2-3,to=1-3, shift left=1pt, no head]
		\arrow[from = 2-3,to=1-3, shift right=1pt, no head]
	\end{tikzcd}
	\]
	The middle map is a homotopic equivalence by Serre \cite{Smale1959}. By comparing the long exact fibrations sequences, the fibers $G$, $SO_2$ must be weakly homotopic equivalent.
\end{proof}
As a consequence of \Cref{EBaer} the group $\Z$, viewed as $\pi_1(\Loop^{eq}(S^2))$, acts as a torsor on the space of varying families of circles $\Gamma_z$ in $S^2 \times [-1,1]$ \emph{with fixed endpoints $\Gamma_0$ and $\Gamma_1$}. Hence, it the view of \Cref{divsetS2iscircle} we can redefine the twisting difference using convex theory.
\begin{defin}
	Let $\xi_1$, $\xi_2$ be two invariant tight contact structures on $S^2 \times [-1,1]$, coinciding near the $z = 0,1$ levels, such that $S^z$ is convex for every $z$. Suppose, moreover, to having fixed two smooth families of dividing circles $\Gamma_z^1$, $\Gamma_z^2$ for the two structures. Then the \emph{convex twisting difference} $\xi_1 - \xi_2$ is the integral difference 
	\[
		[\Gamma^1] - [\Gamma^2] \in \pi_1(\Loop^{eq}(S^2)) \cong \Z\,.
	\]
\end{defin}
Thank to \Cref{parametricdividedetermined} the above definition is independent of the family of dividing circles we choose.

\begin{prop}\label{twdefequal}
	Let $S^2 \times [0,1]$ be an equivariant tomography, and consider two invariant tight contact structures $\xi_1$, $\xi_2$, coinciding at the $0$, $1$ levels. Suppose that there exist two isotopies rel boundary $\phi_t^{(1)}$, $\phi_t^{(2)}$ such that $\xi_i' \coloneqq T\phi_1^{(i)}(\xi_i)$ satisfies the following:
	\begin{enumerate}
		\item the surfaces $S^z$ are $\xi_i$-convex, and
		\item there exists smooth families of dividing circles for $S^z_{\xi_i'}$.
	\end{enumerate}
	Then the following hold:
	\begin{enumerate}
		\item the twisting difference of \Cref{twdif} is independent of the line of fixed point that we choose, and
		\item it coincides with the convex twisting difference.
	\end{enumerate}
\end{prop}
\begin{proof}
	Given a line $F$ of fixed point, let's denote with $d_F(\xi_1, \xi_2)$ the $F$-twisting difference. If $\phi_t$ is an equivariant isotopy rel boundary it must send $F$ to a (possibly equal) line of fixed points. But since the set of such lines is discrete, $F$ must actually be kept fixed. Hence $d_F(\xi_1, \xi_2)$ coincides with $d_F(T\phi_t(\xi_1),\xi_2)$, but also with $d_F(\xi_1, T\phi_t(\xi_2))$. Moreover, thanks to \Cref{parametricdividedetermined} also the convex twisting difference is invariant up to equivariant isotopy rel boundary. Thus, we can suppose that $\xi_i' = \xi_i$.
	
	Every element $\gamma$ of $\Loop^{eq}(S^2)$ determines a tangent vector $\gamma'(a)$ in $T_aS^2$, transversal to $TF$. Thus, a family of such loops gives rise a trivialization of the normal bundle $TF$. On the other hand, if two such families are homotopic, then such identification gives an isotopy of trivializations.
	
	Since the isotopy classes rel boundary of trivializations of $NF$ are isomorphic to $\Z$, there is a \emph{a priori} surjective map 
	\[
		\phi \colon \pi_1(\Loop^{eq}(S^2)) \to \Z\,,
	\]
	that is \emph{a fortiori} a bijection due to \Cref{EBaer}.
	
	We just need to prove that the map $\phi$ is coherent with the torsion action given in \Cref{twdif}, i.e.\ that 
	\[
		\phi([\Gamma^1]- [\Gamma^2]) = d_F(\xi_1, \xi_2)\,.
	\] 
	
	For proving it, let's consider the tangent space $T_pS^z$ as it varies trough $p \in F$. At each point the dividing curve $\Gamma^i$, tangent to $T_pS^z$, is transversal to the characteristic foliation $S^z_{\xi_i}$, and hence to the plane $\xi_i(p)$. Thus the difference $d_F(\xi_1,\xi_2)$, computed looking at the rotation of $\xi_2$ with respect of $\xi_1$, must coincide with the $\phi([\Gamma^1]- [\Gamma^2])$, defined using the tangent vectors of the dividing curves.
\end{proof}


We have hence shown that with the correct assumptions our various definitions are coherent. Instead of proving already that we are in the scenario of \Cref{twdefequal}, we will keep that for later, since it will be already part of the proof of \Cref{thm:mainthmtomographyStwo}.
\subsection{Proof of Theorem \getrefnumber{thm:mainthmtomographyStwo}}



Let's start from an invariant tight contact structure $\xi$ on an equivariant tomography $S \times [-1, 1]$, $S = S^2$.


Supposing $S^{z}$ to be Morse--Smale, and hence convex, in a neighborhoods of $z=0,1$, we want to show how to manipulate the characteristic foliations $S_\xi^z$ by an isotopy rel boundary of $\xi$ until we arrive at a contact structure for which each $S^z$ is a convex surface.

We have seen in \Cref{CinftysurfMorseSmale} that, given an adequate surface $u$-oriented in a contact 3-manifold, by perturbing the contact structure we can arrange the surface to have a characteristic foliation of Morse--Smale type, so that the surface is convex by \Cref{MorseSmaleconvex}. Now, however, we are dealing with a 1-parametric family of surfaces, and the Morse–Smale condition is not generic for 1-parametric families of singular foliations. Nonetheless, there are some simple ``genericity'' assumptions that we may impose. The key to this are the fact that the contact condition is open in the $C^1$-topology, and equivariant Gray stability for invariant contact structures. Therefore, any property that we may impose on a film of singular foliations by a $C^1$-small perturbation can be achieved by a corresponding isotopy of contact structures. 

Let's start with a simple lemma, that allows us to use some extra information given by the contact structure.
\begin{lemma}
	There can be no accumulation point, in terms of the parameter $z$, of retrograde hyperbolic-hyperbolic connections, i.e.\ connections from a negative to a positive hyperbolic point.
\end{lemma}
\begin{proof}
	 Since the statement already holds without any group action, we can ignore the latter for this proof.
	 
	 Arguing by contradiction, assume that we had such an accumulation point at $z = 0$. 
	 
	 Let $z_\nu$, $\nu \in \mathbb{N}$, be a sequence of points $z_\nu$ with $\lim_{\nu \to \infty} z_\nu = 0$ such that each $S_\xi$ contains a retrograde connection, accumulating at such a connection for $\nu \to \infty$. We may choose local coordinates in such a way that we are dealing with a contact structure on $\R^2 \times [-1, 1]$, with a positive (respectively, negative) hyperbolic point at $(\pm1, 0, z_\nu)$, including $z_\infty = 0$, and with the line segment between these to points (for each $\nu$) along the $x$-axis being a retrograde connection. 
	 
	 On each of these slices $S^{z_\nu}$, the vector field defining the characteristic foliation is a positive multiple of $\partial_x$ along the connecting separatrix, and hence $\beta_{z_\nu}$ is a positive multiple of $dy$ (if we take $\Omega = dx \wedge dy$ as area form). From the limiting process it follows that $\dot{\beta}_z|_{z =0}$ is a multiple of $dy$ along the separatrix $\gamma_0$ from $(-1, 0, 0)$ to $(1, 0, 0)$, and thus $(\beta_z \wedge \dot{\beta}_z)|_{z =0}$ = 0 there. 
	 
	 Moreover, there has to be a point on $\gamma_0$ where $w_0 = 0$ and $dw(\partial_x) \ge 0$, since $\gamma_0$ is going from a negative point, where $w_0 < 0$, to a positive one, where $w_0 > 0$. But the contact condition \eqref{eq:contcond} would be negated, since
	 \[
	 	w_0\,d\beta_0 + \beta_0 \wedge (dw_0 - \dot{\beta}_0)(\partial_y, \partial_x) = (\beta_0 \wedge dw_0)(\partial_x, \partial_y)\le 0\,,
	 \]
	 while $\Omega(\partial_x, \partial_y) > 0$.
\end{proof}
\begin{prop}\label{genfilm}
Up to a further small equivariant perturbation of the tomography (equiv.\ a perturbation of the contact structure), we can suppose the following:
\begin{enumerate}
	\item each $S_\xi^z$ contains only finitely many singular points, and --- on finitely many $z$–levels --- at most two degenerate one;
	\item there are no trajectories from a negative to a positive singularity in $S_\xi^z$, except on a finite number of levels $z_1,\dots, z_k$ , where all critical points are non-degenerate and there are exactly two retrograde hyperbolic-hyperbolic connections;
	\item on such levels there are no other connections between hyperbolic points.
\end{enumerate}
\end{prop}
\begin{proof}
	The idea is to readapt the concepts present in \cite[Prop.\ II.2.1]{sotomayorgenericoneparameter}, with the important property of not having closed orbits, nor a poly-cycle as $\omega$- or $\alpha$-limit. By Poincaré-Bendixson Theorem, no other $\alpha$- or $\omega$-limits can appear outside singularities \cite[Thm.\ 4.7.8]{geigesIntroductionContactTopology2008}.
	
	Our conditions are open, we want to show that they are dense. Let's consider the 1-parameter vector field 
	\[
		X \colon S^2 \times [0,1] \to TS^2\,,
	\]
	defining the film of characteristic foliations. By assumption it is a $G$-map, where $G = \Z/2\Z$, with the action on the left-hand side being
	\[
		g\cdot(p,v) = (g\cdot p, - d_pg[v])\,.
	\]
	
	By \Cref{Gtrasv} we know that we can find a new map $X'$, $C^\infty$-close to $X$, such that $X$ is transversal to the zero-section $s_0$. With the same reasoning as in \Cref{nondegequiv}, if we substitute $X'$ with $Z \coloneqq X' \circ h^{-1}$, then $Z$ is a 1-parameter family of vector fields, close to $X$ and transversal to $s_0$.
	
	Next, observe that the zero-locus $S$ of $Z$ is transversal to $S \times \{t_0\}$ at some $(x_0, t_0)$, if and only if $x_0$ is a non-degenerate singular point for $Z(\cdot,t_0)$. The idea then is to work on this condition in order to find a suitable perturbation of $Z$ with the desired properties. 
	
	Let's firstly recall in which situation we are, that is unique for our situation.
	\begin{enumerate}
		\item Since $S$ is a sphere and $\xi$ is tight, $S^z_\xi$ cannot have closed orbits.
		\item Since $S$ is a sphere and $\xi$ is tight, $S^z_\xi$ cannot have poly-cycle as $\alpha$- and $\omega$-limits. In fact, such poly-cycles would be made of singularities of the same sign, and hence would be negated by tightness \cite[Prop.\ 6.4.33]{geigesIntroductionContactTopology2008}.
		\item Since $S^z_\xi$ is a vector field coming from a characteristic foliation, then the divergence at every singular point must be non-zero.
	\end{enumerate}
	
	With that in mind, we just need to observe that the required perturbation is contained in \cite[Prop.\ II.2.1]{sotomayorgenericoneparameter}. After applying the transversality reasoning, that we re-proposed with the correct readaptation, the rest of the perturbation are \emph{local around singular points}. Moreover, the quoted proof does more than we need. In fact, the perturbation Sotomayor calls $\eta^{(1)}$ could have \emph{a priori} non-degenerate non-generic singular points at some levels. Recall: these are traceless singularities transversal to the zero-section. However, the vector field we started with didn't have such singularities, and hence they cannot appear via small perturbations.
	
	Moreover, since $\eta^{(1)}$ is obtained by perturbations around singular points, we can readapt them by carrying them out equivariantly. So we have fulfilled condition $1)$ of our statement.
	
	Finally, since retrograde connections constitute a discrete sets in term of the $z$-parameter, we know that the levels in which they appear must be finite. Hence, we can make small perturbations, analogous to the one used by Peixoto, such that on those levels only two hyperbolic-hyperbolic connections appear, being the retrograde ones, and such that no degenerate singularity manifests.
\end{proof}

For the next lemma recall that an \emph{elimination pair} is a configuration of an hyperbolic and elliptic point, of same sign, connected by a separatrix of the first.
\begin{lemma}
	Let $S$ be an $u$-oriented surface having Morse--Smale characteristic foliation, embedded in an invariant tight contact structure, and let's consider an elimination pair inside $S$. Then we can equivariantly $C^\infty$-perturb $S$ such that the elimination pair disappears.
\end{lemma}
\begin{proof}
	The action $u|_S$ must send the separatrix $(x_e, \gamma, x_h)$ necessarily into one of the same type \emph{and different sign}. So it must be another triple, away from the one we started with.
	
	Since the classical construction \cite[Lemma 4.6.26]{geigesIntroductionContactTopology2008} is local around $\gamma$, we can just apply it in an equivariant way, on $(x_e, \gamma, x_h)$ and $(u(x_e), u(\gamma), u(x_h))$.
\end{proof}
\begin{figure}[t]
	\centering
	\begin{tikzpicture}[x=0.75pt, y=0.75pt,>=Stealth]
		\draw (-50,0) -- (50,0);
		\draw (30:40) -- (210:40);
		\draw (-50,0) -- +(25,20);
		\draw (50,0) -- +(25,20);
		\draw (30:40) -- +(20,20);
		\draw (210:40) -- +(20,20);
		\draw[decoration={markings, mark = at position 0.5 with \arrow{>}},postaction={decorate}] (25,25) ..controls +(-150:0.2cm) and +(60:0.2cm).. (0,0);
		\draw[->] (-70,10) arc [x radius=10, y radius= 5, start angle = 170, end angle = 440];
		
		\draw[->] (-15,15) -- +(-135:17);
		\draw[->] (-15,15) -- +(0:15);
		
		
		\draw[->] (0,0) -- (-45:30);
		\node at (-45:30) [font = \small, shift = {(8,5)}] {$\alpha_p^\dagger$};
		\node at (80,0) [font = \small] {$T_pS$};
		\node at (-45,-25) [font = \small] {$\xi_p$};
		\node at (10,25) [font = \small] {$\vec{\gamma}$}; 
		
		\begin{scope}[shift = {(200,0)}]
			\draw (0,0) circle [radius=40];
			\filldraw (90:40) circle [radius=1.5pt];
			\filldraw (310:40) circle [radius=1.5pt];
			\node at (0,0) [font = \small,shift = {(0,-25)}] {$\theta$};
			\draw (0,5) -- (90:40);
			\draw (310:5) -- (310:40);
			\draw (0,0) circle [radius=5];
			\filldraw (0,0) circle [radius=1];
			\draw[->] (80:15) arc [radius=15, start angle = 80, end angle = 330];
			\node at (30:13) [font = \small] {$\vec{\gamma}$};
		\end{scope}
	\end{tikzpicture}
	\caption{The angle between the oriented planes $T_pS$ and $\xi_p$}
	\label{figangolo}
\end{figure}
\begin{lemma}\label{tomographynoretrograde}
	Let $S$ be an $u$-oriented 2–sphere embedded in an invariant tight contact structure such that $S_\xi$ has finitely many degenerate singular points, and there are no retrograde hyperbolic–hyperbolic connections. Then the graph $G$ in $S$ whose vertices are the positive singular points, and whose edges are the stable separatrices of the positive hyperbolic points, is a connected tree\footnote{The same lemma holds using negative singularities and unstable separatrices.}. Moreover, $S$ is convex, and a dividing set is given by the boundary of a disc-like neighborhood of $G$.
\end{lemma}
\begin{proof}
	Since there are no retrograde saddle-saddle connections, and since the degenerate singular points are isolated, then the proof of \Cref{MorseSmaleconvex} still applies, and we know that $S$ is convex with dividing set given by the boundary of a thickened $G$. In particular, the dividing set is given by a nerve of the annulus between the negative and positive part.
	
	Due to \Cref{divsetS2iscircle}, the graph $G$ is forced to be a connected tree.
\end{proof}
\begin{lemma}\label{nocircleretrograde}
	Let $K$ a topologically trivial Legendrian knot in a tight contact manifold $(M, \xi)$. Let $\Sigma$ a Seifert disk. Then the boundary $\partial \Sigma$ cannot be made of a series of retrograde hyperbolic-hyperbolic connections.
\end{lemma}
\begin{proof}
	We want to show that of a condition implies that the contact framing rotates positively around the surface framing. This will contradicts the Bennequin's inequality $\tb(K) \le -1$ \cite[Thm.\ 4.6.36]{geigesIntroductionContactTopology2008}. 
	
	Consider a retrograde separatrix $\gamma$. Since the Thurston-Bennequin number is indifferent to orientations, we can suppose that the disk is oriented such that is coherent with the orientation of $\gamma$ given by $S_\xi$. At any point $p \in \gamma$, there is a well defined angle $\theta \in S^1$ between $T_pS$ and $\xi_p$, defined as the angle one needs to rotate the \emph{oriented} former plane in order to align it to the latter. See \Cref{figangolo}, where such oriented planes are represented, including the dual vector $\alpha_p^\dagger$.  The rotation is computed positively with respect to $\vec{\gamma}$. At $\gamma(0)$ such angle is $\pi$, at $\gamma(1)$ it is $0 = 2\pi$.
	
	Moreover, at every such point $p \in \gamma$ different from the endpoints we can consider an arbitrary vector $v_p$ such that $\Omega(X_p,v_p) > 0$, where $X$ defines the characteristic foliation. By definition of $X$ we know that $\alpha(v_p) > 0$. Hence $v$ must belong to the \emph{positive component} of $T_pM \setminus \xi_p$ determined by the orientation on $\xi_p$. Consequently the angle at $p$ belongs to $(\pi,2\pi)$.
	
	So, while we are traveling through $\gamma$, the angle between $T_pS$ and $\xi_p$ is contained in $(\pi,2\pi)$. But this implies that such angle is homotopic rel boundary to the linear rotation going from $\pi$ to $2\pi$. But this is a positive half-twist with respect to $T_pS$.

	Consequently, while traversing the full boundary $\partial \Sigma$, the contact framing is doing positive twists around the surface framing.
\end{proof}
\begin{prop}\label{maketomographyallconvex}
	Let $\xi$ be an invariant tight contact structure on a equivariant tomography $S \times [-1,1]$, $S  = S^2$, such that $S^{z}$ are Morse--Smale, and hence convex, near the $z = 0,1$ levels. Then $\xi$ is equivariantly isotopic rel boundary to a contact structure for which each $S^z$, $z \in [-1,1]$ is a convex surface.
\end{prop}
\begin{proof}
	By the discussion above, it suffices to find an isotopy that helps us to get rid of all retrograde saddle–saddle connections. Firstly, we may make all the genericity assumptions discussed up to and with \Cref{genfilm}. Then, by \Cref{tomographynoretrograde}, the surfaces $S^z$ will be convex, except at finitely many levels $z_1, \dots,z_k \in (-1,1)$, where we have two coupled retrograde hyperbolic–hyperbolic connections. For ease of notation, we assume that $z = 0$ is such an exceptional level. We write $x^{\pm}_h$ for the hyperbolic points on that level with a retrograde connection	between them.
	
	Consider the stable separatrix $\gamma$ of $x_h^+$ not coming from $x_h^-$. Its $\alpha$-limit must be a singularity.
	
	If such singularity is an hyperbolic point, then by our genericity assumption such connection must be a retrograde one, coupled with $\gamma$. Hence we are in the situation in which two hyperbolic points, of different signs, are connected through two retrograde connections. However, this situation is ruled out by \Cref{nocircleretrograde}.
	
	Hence $\gamma$ must emanate at a positive elliptic point $x_e^+$. Following the same procedure of \cite[Prop.\ 4.9.7]{geigesIntroductionContactTopology2008}, we observe that we can eliminate the pair $(x_e^+, x_h^+)$ without introducing a new hyperbolic-hyperbolic connections. Moreover, this operation is local on the connected couple, that must be sent ``away'' by the action for sign reasons. Hence we can suppose to work equivariantly.
\end{proof}
This proof shows that by an isotopy rel boundary we may assume that on each slice $S^z$ we have a situation as described in \Cref{tomographynoretrograde}. The last step in the proof of \Cref{thm:mainthmtomographyStwo} is the following lemma.

\begin{lemma}\label{choosezsmoothdividingset}
	Let $\xi$ be a an invariant tight contact structure on an equivariant tomography $S\times[-1,1]$, $S = S^2$, with the properties just described. Then we can choose a dividing set $\Gamma_z$ for	$S^z_\xi$ (consisting of a single circle) that varies smoothly with $z$.
\end{lemma}
\begin{proof}
	The dividing set $\Gamma_z$ can be defined as the boundary of disk-like neighborhoods of the trees $G_z^{\pm}$ defined in \Cref{tomographynoretrograde}, made up of positive or negative singularities and their stable
	or unstable manifolds, respectively, i.e. the collection of flow lines whose $\omega$ or $\alpha$–limit set, respectively, is one of the corresponding singular points. 
	
	As we vary $z$, and pass though a level containing degenerate singular points, the trees $G^{\pm}_z$ can vary their geometries, with new edges and vertices appearing. Remember that $u(G^z_+) = G^z_-$.
	
	Recall \Cref{MorseSmaleconvex}: the dividing set $\Gamma_z$ is given by choosing a nerve of the annulus between the negative and positive parts surrounding $G^\pm_z$ respectively. Thus, if a change in the geometry of $G^\pm_z$ happens at some level $z_k$, we just need define $\Gamma_z$ such that, once the level $z_k$ is reached, the nerve is centered-enough to avoid the change. In this way $\Gamma_{z_k}$ remains a diving circle.
\end{proof}

Now that we have all the pieces, we can prove \Cref{thm:mainthmtomographyStwo}, which we restate.

\setcounter{thmR}{\getrefnumber{thm:mainthmtomographyStwo}}
\addtocounter{thmR}{-1} 
\begin{thmR}
	Let $S^2 \times [-1,1]$ an equivariant tomography. The set of invariant tight contact structure on $S^2 \times [-1,1]$ having fixed characteristic foliation on the boundary, up equivariant isotopy rel boundary, is a $\Z$-torsor.
\end{thmR}
\begin{proof}
	Let's restate the different steps. Let $\xi_1$ and $\xi_2$ be two invariant contact structure as stated, such that $d_F(\xi_1, \xi_2) = 0$ for some line $F$ of fixed point.
	\begin{enumerate}
		\item By \Cref{Girouxbound} we can suppose that, after an equivariant isotopy fixing the boundary, the contact structures coincide near the boundary.
		\item \Cref{Cinftyconvex} allows us to locally find convex 2–spheres $S^2_\pm$ isotopic to $S^2 \times \{\pm1\}$ inside the neighborhoods where the two contact structures coincide. If such new spheres are chosen close enough to the the original ones, the region between is still an equivariant tomography having our convention for the group action. Thus, we can suppose that $S^2 \times \{\pm1\}$ to be already convex.
		\item Thanks to \Cref{genfilm}, up to (different) equivariant isotopies rel boundary, we can suppose that the two tomographies to be standardized.
		\item Thanks to \Cref{maketomographyallconvex} we can use (different) equivariant isotopies rel boundary, to make the tomographies made of convex levels.
		\item Thanks to \Cref{choosezsmoothdividingset} we can suppose to having chosen two smooth families of dividing curves, $\Gamma_z^1$, $\Gamma_z^2$. Hence we can apply \Cref{twdefequal}: the integer $d_F(\xi_1, \xi_2)$ is independent of $F$, and coincides with the convex twisting difference $[\Gamma^1_z] - [\Gamma^2_z]$.
		\item Since $[\Gamma_1]-[\Gamma_2]$ vanishes there is level-preserving equivariant isotopy of $S^2 \times [-1,1]$ such that  $\Phi(\Gamma^1_z)$ coincides with $\Gamma^2_z$. Moreover, since $\xi_0$ and $\xi_1$ coincide near the boundary, we can suppose that such $\Phi$ is the identity in such regions
		\item The contact structures $\Phi^*\xi_1$, $\xi_0$ satisfy the requirement for applying \Cref{uniqueness}. Hence there exist a further equivariant isotopy rel boundary $\Theta$ such that $\Theta^*\Phi^*\xi_1$ coincides with $\xi_0$.
	\end{enumerate}
	This concludes the proof.
\end{proof}
\section{Proof of Theorems \getrefnumber{thmmainthmRthree} and \getrefnumber{thmmainthmDthreecomplementary}}\label{sec:proofmainthmopen}
The aim of this section is to prove \Cref{thmmainthmRthree,thmmainthmDthreecomplementary}. Let's start with the first one, which we restate.

\setcounter{thmR}{\getrefnumber{thmmainthmRthree}}
\addtocounter{thmR}{-1}
\begin{thmR}
	The set of tight contact structures on $\R^3$ up to equivariant isotopy is trivial.
\end{thmR}

The idea is that since $\R^3$ doesn't have a boundary, the twisting difference, that we defined in the previous section, is now meaningless. Let $\xi$ an invariant tight contact structure on $\R^3$.
\begin{prop}\label{Rthreestandardxaxis}
	Suppose that $u$ coincides with the standard rotation, and fix a point $x_0$ on the x-axis. There exists an equivariant isotopy $\phi_t$ of $\R^3$ such that $T\phi_1(\xi)$ is constant on the portion $[x_0,+\infty)$ of the x-axis.
\end{prop}
\begin{proof}
	Let's consider the angle
	\[
		\theta(t) \colon [x_0,+\infty) \to S^1
	\]
	spanned from the $xy$-plane to the contact one at a point $p$, computed positively around the $x$-axis. Let's lift it to a map $\tilde{\theta}$, from $ [x_0,+\infty)$ to the universal cover $\R$.
	
	By denoting the points in $\R^3$ as $\rho\,\hat{v}$, with $\rho$ the norm and $\hat{v} \in S^3$, the desired isotopy $\phi_t$ is then defined as
	\[
		\phi_t(\rho\,\hat{v}) = \rho\,A_{-t\theta(\rho)}\hat{v} \quad \rho \in [x_0, +\infty)\,,
	\]
	with
	\[
		A_{\nu} \coloneqq
		\begin{bmatrix}
			1 & 0 & 0 \\
			0 & \cos\nu & -\sin\nu \\
			0 & \sin\nu & \cos\nu
		\end{bmatrix}\,.
	\]\vspace{0.2cm}
	
	By construction $T\phi_1(\xi)$ has constant value $\ang{\partial_x, \partial_y}$ on the portion $[x_0, +\infty)$ of the $x$-axis.
\end{proof}
\begin{prop}\label{spehrestandardizable}
	Every standard 2-sphere $S \hookrightarrow \R^3$ of any radius in an invariant tight contact structure $\xi$ can be equivariantly perturbed such that it has only two elliptic points: one positive and one negative.
\end{prop}
\begin{proof}
	The non-equivariant statement is contained in \cite[Thm.\ 4.6.30]{geigesIntroductionContactTopology2008}, where the following inequality is proven:
	\begin{equation}\label{exiS}
		\abs{\ang{e(\xi), [S]}} \le 
		\begin{cases}
		0 & S = S^2,\\
		-\chi(S) & \text{otherwise}
		\end{cases}\,.
	\end{equation}
	
	By \Cref{CinftysurfMorseSmale} we can equivariantly perturb $S$ such that $S_\xi$ is Morse--Smale. 
	
	Moreover, from \cite[Prop.\ 4.6.14]{geigesIntroductionContactTopology2008} we know that
	\[
		\ang{e(\xi), [S]} + \chi(S) = 2(e_+ - h_+)\,,
	\]
	where $e_+$ (resp.\ $h_+$) is the number of positive elliptic (resp.\ hyperbolic) points.
	
	From \eqref{exiS} we see that $e(\xi)$ vanishes on $[S]$, and hence
	\begin{equation}\label{eplushplussphere}
		e_+ - h_+ = 1\,.
	\end{equation}
	Thus, we know that there must exists a positive elliptic point $e$. Moreover, from the proof of \cite[Thm.\ 4.6.30]{geigesIntroductionContactTopology2008} we know that only the following two things can happen:
	\begin{enumerate}
		\item $e$ is connected to a positive hyperbolic point, or
		\item the boundary of the basin $D$ of $e$, defined as
		\[
			D = \overline{\{p \to S \mid p \text{ is connected to $e$}\}}\,,
		\]
		consists of a single negative elliptic point.
	\end{enumerate}
	In the first case, the connection must be sent away to itself from the $u$-action, and hence it can be equivariantly eliminated. Since we cannot remain without elliptic points, by Equation \eqref{eplushplussphere}, it must happen that at some point the second case happens, and the basin $D$ is the whole sphere. Hence the foliation is standardized.
\end{proof}
\begin{proof}[Proof of \Cref{thmmainthmRthree}]
	Firstly, let's apply Waldhausen's reduction. Recall the hypothesis we imposed on the involution $u$. Due to Waldhausen \cite{MR236916} such an involution of $S^3$ is necessarily conjugated to an order-2 element in $SO_4$ having an unknot of fixed point. Now is a matter of linear algebra: all these transformation are conjugate to each other. As a consequence all the possible $u$, restricted to $\R^3$, are conjugated one to another. In particular, we can suppose $u$ to be the smooth extension to $S^3$ of
	\[
		(x,y,z) \mapsto (x,-y,-z)\,.
	\]
	
	Let's now consider two contact structure $\xi_1$ and $\xi_2$ as in the statement.
	
	Thanks to \Cref{equalneighorigin} we can find two different equivariant isotopies of $\R^3$ making $\xi_1$ and $\xi_2$ coincide in a small ball $B_\epsilon$ around the origin. Moreover, thanks to \Cref{Rthreestandardxaxis} we can make $\xi_1$ and $\xi_2$ to coincide on the whole $x$-axis, by taking $0 < x_0 <\epsilon$.
	
	From \Cref{spehrestandardizable} we know that there is a local equivariant perturbation $\phi^{n,i}_t$ of the sphere of radius $n$ such that $(\phi^{n,i}_1(S^1_n))_\xi$ is standard. Hence we can substitute $\xi_i$ with $(T\phi_t^{n,i})^{-1}(\xi_i)$, such that the sphere $S^1_n$ has the same standard characteristic foliation, both for $\xi_1$ and $\xi_2$.
	
	Hence, there exists an equivariant contact morphism $\psi$ from $(\NN_1(S^1_n), \xi_1)$ to $(\NN_2(S^1_n), \xi_2)$, that we can read as an isotopy $\psi_t$ from the identity to $\psi$. Again, up to pulling back the contact structure $\xi_2$ through this isotopy, we have reached the situation in which $\xi_1$ and $\xi_2$ coincide to the boundary of spherical shells $I_n$. In such spherical shells can read an equivariant tomography, and hence we can apply the machinery of \Cref{sec:tomography}.
	
	Since $\xi_1$ and $\xi_2$ coincide on the $x$-axis (having the same constant value), the restrictions to $I_n$ have $0$ $F$-torsion difference, with $F$ being the $x$-axis. Then \Cref{thm:mainthmtomographyStwo} tells us that $\xi_1|_{I_n}$ and $\xi_2|_{I_n}$ are equivariantly isotopic rel boundary.
\end{proof}


Let's now move to proving \Cref{thmmainthmDthreecomplementary}. We will follow the same ideas as the previous result. 
\setcounter{thmR}{\getrefnumber{thmmainthmDthreecomplementary}}
\addtocounter{thmR}{-1}
\begin{thmR}
	If such action $u$ preserves the complementary of the standard $3$-ball, then the set of invariant tight contact structures on $(B^3_{\sd})^c$, inducing a fixed characteristic foliation on $S^3 = \partial B^3_{\sd}$, up to equivariant isotopy is trivial.
\end{thmR}
\begin{proof}
	Thanks to Waldhausen's result there exists an infinity-preserving diffeomorphism $W$ of the 3-sphere such that $W \circ u \circ W^{-1}$ is the standard reflection around the $x$-axis. Moreover, $W$ must send the sphere $\partial B_{\sd}$ to a bidimensional compact surface, invariant under reflection, properly containing a portion of the $x$-axis. Let $B$ be the ball $W(B_{\sd})$; we can suppose such ball to contain the origin $\zero$. 
	
	Let $\xi_1$ and $\xi_2$ be two invariant tight structures on the complementary $B^c$ inducing the same characteristic foliation on $\partial B$. By \Cref{Girouxbound} there exists an isotopy rel boundary making them coincide in a small external neighborhood of $B$. Moreover, let's fix a point $x_0$ on the $x$-axis, not belonging to $\partial B$ but contained in such neighborhood. By applying an isotopy that locally coincides with the one given by \Cref{Rthreestandardxaxis}, we can suppose that $\xi_1$ and $\xi_2$ coincide on the portion $[x_0, +\infty)$ of the $x$-axis. 
	
	Finally, following the previous proof, we can fix a sequence of arbitrary large spheres and make $\xi_1$, $\xi_2$ coincide near such spheres. Finally, the regions between such spheres are diffeomorphic to $S^2 \times [-1,1]$, and give rise to an equivariant tomography. The same applies to the region between $\partial B$ and the smaller sphere. Hence we can apply \Cref{thm:mainthmtomographyStwo}.
\end{proof}

\section{Proof of Theorems \getrefnumber{thmmainthmDthree} and \getrefnumber{thmmainthmSthree}}\label{sec:proofmainthmclose}
The aim of this section is to prove \Cref{thmmainthmDthree,thmmainthmSthree}.
\setcounter{thmR}{\getrefnumber{thmmainthmDthree}}
\addtocounter{thmR}{-1}
\begin{thmR}
	If such action $u$ preserves the standard $3$-ball, then the set of invariant tight contact structures on $B^3_{\sd}$, inducing a fixed characteristic foliation on $S^3 = \partial B^3_{\sd}$, up to equivariant isotopy is a $\Z$-torsor.
\end{thmR}

\begin{proof}
	Following the same procedure as in the proof of \Cref{thmmainthmDthreecomplementary}, we can suppose $u$ to be the standard reflection, at the cost of replacing $B^3_{\sd}$ with an arbitrary 3-ball $B$ containing the origin $\zero$.
	
	Given then two contact structures on $B$ inducing the same characteristic foliation on $\partial B$, they can be made to coincide near the boundary by an isotopy rel boundary, thanks to the proof of \Cref{Girouxbound}. Likewise, by \Cref{equalneighorigin} there is a further isotopy that makes the two contact structures coincide near the origin $\zero$, and such that a small standard $2$–sphere $S$ centered at $\zero$ is convex. 
	
	The region $A$ between $S$ and $\partial B$ is diffeomorphic to $S^2 \times [-1,1]$, and equipped with $u$ gives us an equivariant tomography. Now the result follows by applying \Cref{thm:mainthmtomographyStwo} to $A$.
\end{proof}
Let's finish the section with the proof of \Cref{thmmainthmSthree}.
\setcounter{thmR}{\getrefnumber{thmmainthmSthree}}
\addtocounter{thmR}{-1}
\begin{thmR}
	The set of invariant tight contact structures on $S^3$, up to equivariant isotopy, is a $\Z$-torsor.
\end{thmR}
\begin{proof}
	As before, we can suppose $u$ to be the standard rotation.
	
	Given any tight contact structure on $S^3$, we can use \Cref{equalneighorigin} to standardize it near the origin, inside a small 3-ball $B_{\epsilon}$. Now we can apply \Cref{thmmainthmDthree} to the closure of $S^3 \setminus B_{\epsilon}$: via Sch\"{o}nflies theorem such closure can be brought to a standard 3-ball.
\end{proof}
\printbibliography
\end{document}

